\section{Part 11: finite preparation and the properly degenerate local wrapper for P-FIXED}\label{bpP:p11}
\begingroup\sloppy

This part prepares the fixed-Sun Hamiltonian of Theorem~\ref{bpP:p0:fixed} (P-FIXED, general $n\ge2$) for the maximal-case analytic KAM theorem of Part~11a, and transports the resulting tori back to Cartesian coordinates. It follows CP09 \S2.3 (finite-order averaging, polar normalisation of the secular oscillators, the amount of non-degeneracy of the prepared frequency map, Lemma~8 and its polynomial budget) for the properly degenerate Hamiltonian $H=h(I)+\varepsilon f$ of CP09 Theorem~4 (here: $d=p=n$ fast angles and secular oscillators, so $N=d+p=2n$), with every quantitative gap of the source closed (reviews \texttt{p11-route-study.md}, \texttt{p11a-russmann-check.md}: R-CP-1, R-CP-2, A2, A4--A7). The route and its deviations in one paragraph.

\paragraph{Route.} (A) Part~7 (product chart), Part~8 (secular function), Part~9 (diagonalising rotation $\varrho$) put $H_{\rm fix}$ into the form $\tilde h(J)+\varepsilon F(\theta,J,x,y)$ with $F$ holomorphic on a strip and a complex action polydisc, $\theta\in\T^n$ the fast angles (period one), $(x,y)$ the rotated secular Cartesian pairs (Prop.~\ref{bpP:p11:holo-chart}). (B) A base point $J_0$ with Diophantine joint frequency $(\omega,\tilde\Omega)(J_0)$ is chosen by Part~11a-1 (Lemma~\ref{bpP:p11:select}; this replaces CP09 Lemma~3 and the erroneous set $K\cap D_{\gamma_0,\tau_0}$ of CP09 Lemma~3, errata item E-P-CP09-4 in \S\ref{bpP:p11:sec-notes}). (C) \emph{Fast averaging} (CP09 Lemma~4): Part~2's oscillator averaging removes the $\theta$-dependence up to $(\varepsilon M_F)^3$ on the action polydisc of radius $\asymp(\log\frac1\varepsilon)^{-(\tau_0+1)}$ (Prop.~\ref{bpP:p11:avg-fast}). (D) \emph{Polar normalisation and slow averaging}: on the shell $\abs{x_j}^2+\abs{y_j}^2\approx\varepsilon$ the oscillators are put in action--angle form $(\tilde\zeta,\tilde\rho)$, and a second two-scale averaging, now over all $2n$ angles with divisors $\alpha_1$ (fast) and $\varepsilon\alpha_2$ (slow), makes the remainder $O(\varepsilon^{\nu_2+3/2})$ for an arbitrary fixed $\nu_2$ (Lemma~\ref{bpP:p11:aniso-avg}, Prop.~\ref{bpP:p11:avg-slow}). \emph{This replaces CP09 Lemmas 5--7} (equilibrium shift, quadratic diagonalisation, finite Birkhoff normal form): the Birkhoff step is not needed for the Rüssmann gate because only the $J$-derivatives of the prepared frequency map enter the amount of non-degeneracy (Lemma~\ref{bpP:p11:amount}), and the exponentially small slow remainder replaces the $O(\abs{(u,v)}^{\nu_2+1})$ remainder. It also bypasses the errata A5 (polar domain) and A6 (Birkhoff formula) of the review; the price is that $\tilde\Omega(J_0)$ must be Diophantine, which Lemma~\ref{bpP:p11:select} provides for free (deviation D-11-1, \S\ref{bpP:p11:sec-notes}). (E) The \emph{prepared system} $\hat H_\varepsilon=N_\varepsilon+P_\varepsilon$ in rescaled actions $J=J_0+\varepsilon r$, $\tilde\rho=\varepsilon(2\pi+\rho')$ lives on a fixed compact $K_\varepsilon=[-\frac12,\frac12]^{2n}$ with $\vartheta=\frac14$ and $\varepsilon^*=\frac12$, so that \emph{every chart constant of Part~11a-5 is independent of $\varepsilon$ except the amount} $\beta_\varepsilon\asymp\varepsilon^{\bar\mu+2}$ (Prop.~\ref{bpP:p11:prepared}, Lemma~\ref{bpP:p11:amount}); this is the repaired form of CP09 Prop.~2 (R-CP-2) and of CP09 (42)--(47) (A4). (F) Theorem~\ref{bpP:p11:wrapper} applies the KAM input (Definition~\ref{bpP:p11:kam-input}) and transports the tori through the total chart to the tube $U_{\zeta_*}$.

\subsection{Conventions, external facts, and the KAM input}\label{bpP:p11:sec-conv}

\paragraph{Conventions.} Frozen conventions of Part~0: angles of period one, characters $\chi_k=e^{2\pi ik\cdot\theta}$, max norms $\abs\cdot_\infty$ on $\R^m,\C^m$ and induced operator norms. $D(c,r)=\{v\in\C^m:\abs{v-c}_\infty<r\}$ (complex polydisc), $\T_\sigma=\{\theta\in\C^m:\abs{\operatorname{Im}\theta}_\infty<\sigma\}$. For $V\subset\C^m$ open and a function $u$ on $\T_\sigma\times V$, $1$-periodic in the angles, the \emph{weighted norm} is that of \Nref{bpN:p4:norm},
\[\wn u V\sigma:=\sup_{x\in V}\sum_{k}\abs{\hat u_k(x)}\,e^{2\pi\abs k_1\sigma},\]
(Fourier expansion in the angles only; the remaining variables, including oscillator coordinates, play the r\^ole of the parameter $x$); $\sup\abs u\le\wn uV\sigma$ and $\wn uV\sigma\le\coth(\pi\sigma')^{m}\sup_{\T_{\sigma+\sigma'}\times V}\abs u$ (\Nref{bpN:p4:sup-weighted}). $c,c_1,\dots$ denote positive constants; a constant is always listed with the data it depends on. Period one is used for \emph{all} angles of the prepared problem. The period conversion is the canonical chart $\mathcal T(\theta,J,\zeta,\tilde\rho)=(2\pi\theta,J/2\pi,2\pi\zeta,\tilde\rho/2\pi)$ of Remark~\ref{bpP:p2:use5}(3) (a $1$-chart, $dJ\wedge d\theta=d\Lambda\wedge d\lambda$): the Hamiltonian is the same function, frequencies are divided by $2\pi$, and the frozen frequency is $w=\omega_{\rm CP}/2\pi$ (Lemma~\ref{bpP:p1:period-convert}).

\paragraph{External facts.} Each is a standard statement, labelled for reference.
\begin{enumerate}[label=(e\arabic*),ref=(e\arabic*)]
\item\label{bpP:p11:e1} A real analytic function on an open subset of $\R^m$ extends to a holomorphic function on a complex neighbourhood of every compact subset (\Kref{bp:p1:holomorphic-analytic}); the extension is real at real points, and identities between real analytic functions (such as $\varrho^\top\varrho=1$) persist on the connected complex neighbourhood (identity theorem).
\item\label{bpP:p11:e2} Cauchy estimates: if $u$ is holomorphic and $\lvert u\rvert\le S$ on a polydisc $D(c,\delta)\subset\C^m$ (max norm), then $\lvert\partial^\alpha u(c)\rvert\le\alpha!\,S\,\delta^{-\lvert\alpha\rvert}$; the Taylor coefficients of $u$ at $c$ satisfy $\lvert u_\alpha\rvert\le S\delta^{-\lvert\alpha\rvert}$ and $\sum_\alpha\lvert u_\alpha\rvert\delta'^{\lvert\alpha\rvert}\le S(1-\delta'/\delta)^{-m}$ for $\delta'<\delta$ (\Nref{bpN:p4:general-cauchy}).
\item\label{bpP:p11:e3} A continuous function on a compact set is uniformly continuous (Heine--Cantor).
\item\label{bpP:p11:e4} For $\varepsilon\in(0,e^{-1}]$: $\log\frac1\varepsilon\ge1$, and $\varepsilon^{a}(\log\frac1\varepsilon)^b\to0$ as $\varepsilon\to0$ for $a>0$, $b\ge0$.
\item\label{bpP:p11:e5} Termwise integration of a uniformly convergent series and $\int_0^1\sin^a(2\pi t)\cos^b(2\pi t)\,dt=0$ unless $a$ and $b$ are both even (parity of the integrand under $t\mapsto t+\frac12$ and $t\mapsto-t$), with absolute value $\le1$.
\end{enumerate}
Below, (e1)--(e5) refer to these items.

\begin{definition}[Quantitative KAM input, in the form in which it is used]\label{bpP:p11:kam-input}
This is a \emph{citation with exact interface}, not a hypothesis of any lemma of this part, and it is used only in the proof of Theorem~\ref{bpP:p11:wrapper}. Items (K1)--(K4) are \Pref{bpP:p11a:weak-kam-quantitative} with \Pref{bpP:p11a:limit-torus}, \Pref{bpP:p11a:displacement}, \Pref{bpP:p11a:evaluation-bilipschitz} (Files 11a-4, 11a-5). (K5) is proved: it is Corollary~\ref{bpP:p11a:adm-budget-beta} (equivalently Proposition~\ref{bpP:p11a:adm-budget}). (K6) is the one piece that is \emph{not} supplied by File~11a-3; it is a visible dependency on 11a-4/11a-5. Fix $N=2n$ and the squared convention $\rho=2$ of Definition~\ref{bpP:p11a:index-def}. A \emph{KAM datum} is $\mathfrak d=(\mu,\beta,K,\vartheta,s,M,M_0,C_1,\tau_D,\varepsilon^*)$ with $\mu\in\N_{\ge1}$, $\vartheta,s\in(0,1]$, $C_1\ge1$, $M_0>0$, $K\subset\R^N$ compact with $d=\operatorname{diam}_\infty K$, $\varepsilon^*\in(0,\operatorname{meas}K)$, $0<\beta\le\min\{1,M\}$, $\tau_D>\max\{N\mu/2-1,\ N-1\}$. A Hamiltonian $H=h(r)+f(\theta,r)$ on $\T_s\times U$, $U=K+2\vartheta\subset\C^N$ (symplectic form $dr\wedge d\theta$), is \emph{$\mathfrak d$-admissible} (Assumption~\ref{bpP:p11a:sel-data}) if $h,f$ are holomorphic, $f$ is $1$-periodic in $\theta$, both are real on real points, and, with $\omega=\nabla h$,
\begin{enumerate}[label=(A\arabic*)]
\item $\beta_{2,\mu}(\omega,K)\ge\beta$ (Definition~\ref{bpP:p11a:index-def}); \item $\lvert D^kg_2[\omega](c,y)\rvert_{\mathrm{op}}\le M$ for $c\in S^{N-1}$, $1\le k\le\mu+1$, $y\in K_\vartheta$, with $M\ge\beta$; \item $\lVert D\omega\rVert=\lvert D^2h\rvert_{\infty\to\infty}\le C_1$ on $U$; \item $\lVert\omega\rVert^\mu_{\mathrm{dg},K_\vartheta}\le M_0$.
\end{enumerate}
\emph{Input.} For every $\mathfrak d$ there are $\gamma>0$ and $\varepsilon_{\rm adm}>0$, depending only on $\mathfrak d$ (hence on $\beta$ through $\mathfrak d$), such that for every $\mathfrak d$-admissible $H$ with $\sup_{\T_s\times U}\lvert f\rvert\le\varepsilon_f\le\varepsilon_{\rm adm}$:
\begin{enumerate}[label=(K\arabic*)]
\item a compact $\mathcal H\subset K$ with $\operatorname{meas}\mathcal H\ge\operatorname{meas}K-\varepsilon^*$;
\item for $b\in\mathcal H$ a real analytic $\mathcal K_b:\R^N\to\R^N\times\R^N$ (positions, actions) with $\mathcal K_b(\theta+k)=\mathcal K_b(\theta)+(k,0)$ ($k\in\Z^N$) --- equivalently $u_b:=\mathcal K_b-(\mathrm{id},b)$ is $\Z^N$-periodic --- $D\mathcal K_b$ injective, $\mathcal K_b^*(dr\wedge d\theta)=0$, $\mathcal K_b(\theta)=\mathcal K_b(\theta')$ only if $\theta'-\theta\in\Z^N$, and a frequency $w_b\in\R^N$ that is $(\gamma,\tau_D)$-Diophantine with $s\mapsto\mathcal K_b(\theta+sw_b)$ having derivative $X_H(\mathcal K_b(\theta+tw_b))$ at every $t$;
\item $\sup_{b,\theta}\lvert u_b(\theta)\rvert_\infty\le\Delta$ for a number $\Delta=\Delta(\varepsilon_f;\mathfrak d)$;
\item the evaluation map $(b,\theta)\mapsto\mathcal K_b(\theta)$ is injective on $\mathcal H\times[0,1)^N$, its image $\mathcal B$ is Borel, and $\operatorname{meas}\mathcal B\ge2^{-2N}\operatorname{meas}\mathcal H$;
\item \emph{(lower budget; proved)} if all entries of $\mathfrak d$ except $\beta$ are held fixed, then there is $c_*>0$ depending only on them with $\varepsilon_{\rm adm}\ge c_*\beta^{B_0}$ for $0<\beta\le\min\{1,M\}$, where $B_0=\mu(\mu+1)+2$ (Corollary~\ref{bpP:p11a:adm-budget-beta}, $\rho=2$; the threshold actually used is the monotone minorant $\varepsilon^\flat_{\mathrm{adm}}\le\varepsilon_{\mathrm{adm}}$ of File 11a-5, Proposition~\ref{bpP:p11a:budget}(d), which gives $\varepsilon^\flat_{\mathrm{adm}}\ge c_*\beta^{B_0}$ and part (a) of which makes it valid for any actual $M_0'\le M_0$; since the actual values of $M,M_0,C_1$ and $d$ of a family only need to be bounded above, one applies Proposition~\ref{bpP:p11a:adm-budget} with $a_\beta=1$ and all other rates $0$, which gives the same $B_0$);
\item \emph{(displacement power; dependency on 11a-4/11a-5, not on 11a-3)} with the same entries fixed there are $c_D>0$, $\kappa_D>0$, $a_D\ge0$ with $\Delta\le c_D\,\varepsilon_f^{\kappa_D}\beta^{-a_D}$ whenever $\varepsilon_f\le\varepsilon_{\rm adm}$ and $\beta\in(0,\min\{1,M\}]$.
\end{enumerate}
\end{definition}

\subsection{Step A: the Hamiltonian in working coordinates}\label{bpP:p11:sec-A}

\paragraph{The fixed data and the patch.} $n\ge2$, $G,m_0>0$, $\kappa=Gm_0$, $m\in(0,\infty)^n$, $a\in\R^n$ with $0<a_1<\dots<a_n$, $\eta>0$ (Theorem~\ref{bpP:p0:fixed}; Lean index shift as in Part~12). $\zeta_*,d,D$ are those of Definition~\ref{bpP:p6:tube-constants}; $\Lambda^*_i=m_i\sqrt{\kappa a_i}$, $\rho_*=\rho_*(m,\kappa,a)$ of Proposition~\ref{bpP:p8:perturbation-analytic}. Let $\delta_*>0$ be the constant of Lemma~\ref{bpP:p12:forward-tube} for $\zeta=\zeta_*$ (the proof of that lemma uses $n=2$ only through the number of planets; Lemma~\ref{bpP:p11:forward-tube} restates it). Choose, by Theorem~\ref{bpP:p10:interface} (with the prescribed neighbourhood of $a$ so small that $\abs{\Lambda(a')-\Lambda^*}_\infty<\min\{\rho_*,\delta_*\}/8$), the patch $a_0,r_0,K_a$ and put $K_\Lambda=\Lambda(K_a)$, the index $\bar\mu=\mu_0\ge1$ and amount $\bar\beta>0$ of $\tilde\varpi=(\nu,\Omega)$ in the squared convention (Theorem~\ref{bpP:p10:interface}(d)), depending only on the data. In the period-one working variable $J=2\pi\Lambda$ put
\[\begin{gathered}\tilde h(J)=h(J/2\pi),\quad\omega(J)=\nabla\tilde h(J)=\tfrac1{2\pi}\nu(J/2\pi),\quad\tilde\Omega(J)=\tfrac1{2\pi}\Omega(J/2\pi),\\ \varpi_J=(\omega,\tilde\Omega):\ 2\pi B^\Lambda\to\R^{2n},\qquad K_J=2\pi K_\Lambda,\end{gathered}\]
with $\Omega_j=2\sigma_j$ (Remark~\ref{bpP:p9:dict}) and $\tilde\Omega_j<0$.

\begin{lemma}[Index and amount in the working variable]\label{bpP:p11:norm-dict}
$\varpi_J$ is real analytic and R-nondegenerate on $2\pi B^\Lambda$, and $\beta_{2,\bar\mu}(\varpi_J,K_J)\ge(2\pi)^{-2-\bar\mu}\bar\beta=:\beta_J>0$.
\end{lemma}
\begin{proof}
\proofstep{Constant factor and dilation}\label{stepP:p11:nd} $\varpi_J(J)=(2\pi)^{-1}\tilde\varpi(J/2\pi)$. Hence $g_2[\varpi_J](c,J)=(2\pi)^{-2}g_2[\tilde\varpi](c,J/2\pi)$ (Definition~\ref{bpP:p11a:index-def}) and $D_J^\nu=(2\pi)^{-\nu}D_\Lambda^\nu$, so $\max_{\nu\le\bar\mu}\lvert D^\nu_Jg_2[\varpi_J](c,J)\rvert\ge(2\pi)^{-2-\bar\mu}\max_{\nu\le\bar\mu}\lvert D^\nu_\Lambda g_2[\tilde\varpi](c,\Lambda)\rvert\ge(2\pi)^{-2-\bar\mu}\bar\beta$ for $c\in S^{2n-1}$, $J\in K_J$ (Theorem~\ref{bpP:p10:interface}(d)). R-nondegeneracy is invariant under a positive diagonal rescaling and an affine change of the variable (Theorem~\ref{bpP:p10:rnd}(c), Remark~\ref{bpP:p10:scaling}).
\end{proof}

\begin{lemma}[Base point]\label{bpP:p11:select}
Put $\tau_0=\max\{n\bar\mu,\,2n\}$ and $\varepsilon_0^*=\tfrac12\operatorname{meas}K_J$. There are $\gamma_0\in(0,1]$ and $J_0\in\operatorname{int}K_J$, depending only on the data, with $\abs{k\cdot\varpi_J(J_0)}\ge\gamma_0\abs k_1^{-\tau_0}$ for all $k\in\Z^{2n}\setminus\{0\}$. In particular $\omega(J_0)$ and $\tilde\Omega(J_0)$ are $(\gamma_0,\tau_0)$-Diophantine in $\Z^n$ (Definition~\ref{bpP:p0:diophantine}).
\end{lemma}
\begin{proof}
\proofstep{Hypotheses of the measure estimate}\label{stepP:p11:sel-hyp} Apply Corollary~\ref{bpP:p11a:power-law} with $N=n$, $m=2n$, $\rho=2$, $\mu=\bar\mu$, $\omega=\varpi_J$, $\tilde\omega_k=\varpi_J$ for all $k$, the compact set $K_J$ (nonempty, positive measure, $d=\operatorname{diam}_\infty K_J>0$), $\vartheta>0$ so small that $(K_J)_\vartheta\subset 2\pi B^\Lambda$, $\beta=\beta_J$ (Lemma~\ref{bpP:p11:norm-dict}) and $M=\max\{\beta_J,M_1\}$ of Lemma~\ref{bpP:p11a:M-bound}. (D1) is Lemma~\ref{bpP:p11:norm-dict}, (D2) is Lemma~\ref{bpP:p11a:M-bound}, (D3) is trivial because $\tilde\omega_k=\omega$. The exponent condition is $\tau_0>m\mu/\rho-1=n\bar\mu-1$.
\proofstep{Choice of $\gamma_0$}\label{stepP:p11:sel-gam} By the Corollary $\operatorname{meas}(K_J\setminus H_\gamma)\le C_{\rm Dio}\gamma^{2/\bar\mu}$ for $\gamma\le\gamma_*(\varepsilon^*_0)$; take $\gamma_0=\min\{1,\gamma_*(\varepsilon_0^*)\}$. Then $\operatorname{meas}H_{\gamma_0}\ge\operatorname{meas}K_J-\varepsilon_0^*>0$. Here $H_{\gamma_0}=\{b\in K_J:\abs{k\cdot\varpi_J(b)}\ge\gamma_0\abs k_1^{-\tau_0}\ \forall k\ne0\}$ since the map does not depend on $k$.
\proofstep{Interior point}\label{stepP:p11:sel-int} $K_J\setminus\operatorname{int}K_J$ is the boundary of a \emph{rectangle} (the map $a\mapsto2\pi\Lambda(a)$ acts coordinatewise and is increasing, so it maps the cube $K_a$ to a product of intervals, not in general a cube), which is a null set. Hence $H_{\gamma_0}\cap\operatorname{int}K_J\neq\emptyset$; take $J_0$ in it. The restriction to $k=(k_1,0)$ and $k=(0,k_2)$ is the last sentence.
\end{proof}
This replaces CP09 Lemma 3 (\texttt{p11a-russmann-check}, A7): the object is the set of \emph{parameters} $b$ with $\varpi_J(b)\in\mathcal D_{\gamma_0,\tau_0}$, not $K\cap\mathcal D$; the exponent of CP09 ($\tau_0\ge d\mu_0$) is replaced by the weaker $\tau_0>n\bar\mu-1$ proved in File~11a-1 (deviation D5 there).

\begin{proposition}[Holomorphic secular form near the base point]\label{bpP:p11:holo-chart}
Let $J_0,\gamma_0,\tau_0$ be as in Lemma~\ref{bpP:p11:select}. There are $r_C\in(0,1]$, $\sigma_F\in(0,\frac1{4}]$, $r_z\in(0,1]$, $M_F,L,B_\varrho,B'_\varrho\ge1$ and a function $F$, all depending only on the data (and $J_0$), such that, with $\mathcal J=D(J_0,2r_C)\subset\C^n$ and $z=(x,y)\in\C^{2n}$ (position $x$, momentum $y$):
\begin{enumerate}[label=(\alph*)]
\item $F$ is holomorphic on $\T_{2\sigma_F}\times\mathcal J\times B_{r_z}$ ($B_{r_z}=D(0,r_z)\subset\C^{2n}$), $1$-periodic in $\theta$, real on real points, and $\abs F\le M_F/\coth(\pi\sigma_F)^n$ there; hence $\wn F{\mathcal J\times B_{r_z}}{\sigma_F}\le M_F$ in the sense of \S\ref{bpP:p11:sec-conv} with $z$ as a parameter. Moreover $\omega,\tilde\Omega$ are holomorphic on $\mathcal J$ and $\lvert D^j\varpi_J\rvert_{\mathrm{op}}\le L$ for $0\le j\le\bar\mu+1$ there (so $\lvert D\omega\rvert,\lvert D\tilde\Omega\rvert,\lvert\varpi_J\rvert\le L$).
\item (\emph{identity}) $\Psi_c:=\Psi^{(n)}\circ\Phi_\varrho^{-1}\circ\mathcal T_1$, with $\mathcal T_1(\theta,J,x,y)=(2\pi\theta,J/2\pi,x,y)$ and $\varrho:D(\Lambda_0,r_C/\pi)\to O(n,\C)$, $\Lambda_0=J_0/2\pi$, the holomorphic extension of the analytic rotation of Proposition~\ref{bpP:p9:local-analytic} (diagonalising $Q_h$; $\varrho^\top\varrho=1$ on the whole polydisc, $\varrho$ real with values in $SO(n)$ at real points, $\lvert\varrho\rvert_{\infty\to\infty}\le B_\varrho$ and $\lvert\varrho^\top\partial_{\Lambda_l}\varrho\rvert_{\infty\to\infty}\le B'_\varrho$ for all $l$ there), is a symplectic $1$-chart on $\mathcal O_c=\{(\theta,J,z)\in\R^{2n}\times\R^{2n}:J\in\mathcal J\cap\R^n,\ \abs z_\infty<r_z\}$ with values in the collision-free domain, and for all real $\varepsilon$ and $(\theta,J,z)\in\mathcal O_c$: $H_{\rm fix}\bigl(\Psi_c(\theta,J,z)\bigr)=\tilde h(J)+\varepsilon F(\theta,J,z)$. The symplectic form on the working coordinates is $dJ\wedge d\theta+dy\wedge dx$ (momenta wedge positions, the frozen order $\Omega=\sum dz_2\wedge dz_1$ of Part~0; the positions are $(\theta,x)$, the momenta $(J,y)$), and \emph{all} charts below are symplectic for this form.
\item (\emph{lattice}) $\Psi_c$ satisfies (Q1)--(Q2) of Theorem~\ref{bpP:p1:torus-transport} for the lattice $\iota_k=((k,0),0)$ shifting $\theta$ and $\sigma_k=0$ (the secular positions $x$ are not shifted; the secular lattice enters after the polar chart, Lemma~\ref{bpP:p11:total-chart}).
\item (\emph{mean}) $F_0:=\hat F_0$ (the Fourier average over $\theta$) is holomorphic on $\mathcal J\times B_{r_z}$ and
\[F_0(J,z)=C(J)+\sum_{j=1}^n\Omega_j(J)\tfrac12(x_j^2+y_j^2)+R_4(J,z),\qquad\abs{R_4(J,z)}\le M_F\bigl(\abs z_\infty/r_z\bigr)^4,\]
with $C(J)=C_0(a(J/2\pi))$, $\Omega_j(J)=2\pi\tilde\Omega_j(J)$ and real analytic $C$. Both models would give this form (Corollary~\ref{bpP:p8:same-jets}); only $F_0$ of the fixed-Sun model is used.
\end{enumerate}
\end{proposition}
\begin{proof}
\proofstep{Rotation and its extension}\label{stepP:p11:hc-rot} $Q_h$ has simple spectrum on the ball $\mathcal W\ni\Lambda_0$ (Theorem~\ref{bpP:p10:interface}(a), $J_0\in K_J$), so $\sigma_j,\varrho$ are real analytic on a ball $D_\R(\Lambda_0,t_1)$ (Proposition~\ref{bpP:p9:local-analytic}). By (e1) $\varrho$ extends holomorphically to a polydisc $D(\Lambda_0,t_2)$, $t_2\le t_1$; the identities $\varrho^\top\varrho=1$ and $\det\varrho=1$ hold on the real ball, hence on $D(\Lambda_0,t_2)$ (identity theorem), so $\varrho$ takes values in $O(n,\C)$; at real points it takes values in $SO(n)$, which are Euclidean isometries. \emph{Off the real locus complex orthogonality gives no isometry} (the rotation of angle $it$ has $\lvert Re_1\rvert_\infty=\cosh t>1$), and the derivative of $\varrho$ at $\Lambda_0$ may be arbitrarily large; so no normalisation such as $\lvert\varrho\rvert\le2$ is claimed. Instead fix $r_C\le\min\{\pi t_2/2,\ \pi\rho_*/8\}$ and put $B_\varrho=\sup\lvert\varrho\rvert_{\infty\to\infty}$, $B'_\varrho=\max_l\sup\lvert\varrho^\top\partial_{\Lambda_l}\varrho\rvert_{\infty\to\infty}$ on $D(\Lambda_0,r_C/\pi)$ (finite: continuous on a compact subset of $D(\Lambda_0,t_2)$), enlarged to $\ge1$.
\proofstep{The map $\Phi_\varrho^{-1}$ on the complex domain}\label{stepP:p11:hc-lift} By Proposition~\ref{bpP:p9:lift}, $\Phi_\varrho^{-1}(\tilde\lambda,\Lambda,\tilde\xi,\tilde\eta)=(\tilde\lambda+H_\varrho,\Lambda,\varrho\tilde\xi,\varrho\tilde\eta)$ with $H_{\varrho,l}=\tilde\eta^\top(\varrho^\top\partial_{\Lambda_l}\varrho)\tilde\xi$; the identity of the lift is polynomial in the entries of $\varrho,\partial\varrho$, hence holds on the complex domain, and $H_\varrho$ is holomorphic there. For $\lvert z\rvert_\infty<r_z$ (with $\tilde\xi=x$, $\tilde\eta=y$): $\lvert\varrho x\rvert_\infty,\lvert\varrho y\rvert_\infty\le B_\varrho r_z$ and $\lvert H_\varrho\rvert_\infty\le\lvert y\rvert_1\,B'_\varrho\lvert x\rvert_\infty\le nB'_\varrho r_z^2$. Require
\[r_z\le\min\Bigl\{1,\ \frac{\rho_*}{8B_\varrho},\ \Bigl(\frac{\pi\sigma'}{2nB'_\varrho}\Bigr)^{1/2}\Bigr\}\]
($\sigma'$ of the next step), so that the rotated secular variables lie in $\mathcal N(\rho_*/8)$-range and the angle shift has modulus $\le\pi\sigma'/2$ in $\lambda$ (that is $\sigma'/4$ in $\theta$ units).
\proofstep{Strip extension of $\hat f_{\rm fix}$}\label{stepP:p11:hc-strip} Part~8 proves the holomorphy of $\hat f_{\rm fix}$ in $w'=(\xi,\Lambda,\eta)\in\mathcal N(\rho_*)$ only for real $\lambda$ (Proposition~\ref{bpP:p8:perturbation-analytic}(b)). Here the angular strip is needed. By Proposition~\ref{bpP:p7:holomorphic} $s_{ij}(\lambda,w')$ is holomorphic on $\{\abs{\operatorname{Im}\lambda}<\frac14\}\times\mathcal N(\rho_0)$ and $2\pi$-periodic; by Proposition~\ref{bpP:p8:perturbation-analytic}(a) $\operatorname{Re}s_{ij}\ge\frac12(a_j-a_i)^2$ at real $\lambda$ and $w'\in\mathcal N(\rho_*)$ (the $a_i$ here are the ordered axes of the data, $a_i^*$ of Part~8). By (e3) applied to $[0,2\pi]^n\times\overline{\mathcal N(\rho_*/2)}$ (and the $2\pi$-periodicity in $\operatorname{Re}\lambda$) there is $\sigma'\in(0,\frac14]$ with $\operatorname{Re}s_{ij}\ge\frac14(a_j-a_i)^2$ for $\abs{\operatorname{Im}\lambda}_\infty\le2\pi\sigma'$ (a strip of period-one width $\sigma'$ in $\theta=\lambda/2\pi$) and $w'\in\mathcal N(\rho_*/2)$. On this set $\hat f_{\rm fix}=-G\sum_{i<j}m_im_js_{ij}^{-1/2}$ (principal root) is holomorphic and bounded by $c(m,G,a)$.
\proofstep{Definition of $F$ and (a)}\label{stepP:p11:hc-F} $F:=\hat f_{\rm fix}\circ\Phi_\varrho^{-1}\circ\mathcal T_1$ is holomorphic where the composition is defined: $\theta\in\T_{2\sigma_F}$ with $2\sigma_F\le\sigma'/2$, the shift $\abs{\operatorname{Im}H_\varrho}\le nB'_\varrho r_z^2\le\pi\sigma'/2$ (Step~\ref{stepP:p11:hc-lift}), and $\Lambda\in D(\Lambda_0,r_C/\pi)$ with $\lvert\Lambda_0-\Lambda^*\rvert_\infty<\rho_*/8$, $r_C/\pi\le\rho_*/8$, so $\lvert\Lambda-\Lambda^*\rvert<\rho_*/4$; $\lvert\varrho x\rvert_\infty<\rho_*/8$. Here $\sigma'$ is the strip width of Step~\ref{stepP:p11:hc-strip}, which depends only on $\rho_*$ and the data and not on $r_z$ (so the choices are not circular). $F$ is $1$-periodic in $\theta$ ($\hat f_{\rm fix}$ is $2\pi$-periodic in $\lambda$ and $H_\varrho$ does not depend on $\tilde\lambda$), real at real points ($\varrho$, $\hat f$ real). Its sup is some $M_\infty$; put $M_F=\max\{1,\coth(\pi\sigma_F)^nM_\infty\}$ and use \Nref{bpN:p4:sup-weighted}(b) with the extra strip $\sigma_F$. $L$ bounds $\varpi_J$ and its derivatives up to order $\bar\mu+1$ on $\mathcal J$ (holomorphic on a neighbourhood of $\overline{\mathcal J}$ after shrinking $r_C$; Cauchy, (e2)).
\proofstep{(b) Identity and chart}\label{stepP:p11:hc-id} For real points the composition $\Psi^{(n)}\circ\Phi_\varrho^{-1}\circ\mathcal T_1$ lies in $\Psi^{(n)}(\mathcal N_\R(\rho_*)\times\R^n)\subset W$ (Proposition~\ref{bpP:p8:perturbation-analytic}(a)); Theorem~\ref{bpP:p7:product-chart}(c) gives $H_{\rm fix}\circ\Psi^{(n)}=h(\Lambda)+\varepsilon\hat f_{\rm fix}$ and $\Phi_\varrho$ preserves $\Lambda$ (Proposition~\ref{bpP:p9:lift}), $\mathcal T_1$ maps $J\mapsto\Lambda=J/2\pi$: $h(\Lambda)=\tilde h(J)$. Symplecticity: $\Psi^{(n)}$ (Theorem~\ref{bpP:p7:product-chart}(a)), $\Phi_\varrho^{-1}$ (Proposition~\ref{bpP:p9:lift}), $\mathcal T_1$ ($d\Lambda\wedge d\lambda=dJ\wedge d\theta$), Remark~\ref{bpP:p1:typical}(3); the chart is real analytic because $F$ and the maps are (Lemma~\ref{bpP:p1:hol-chart}).
\proofstep{(c) Lattice}\label{stepP:p11:hc-lat} $\mathcal T_1(\theta+k,\cdot)=\mathcal T_1(\theta,\cdot)+(2\pi k,0,0,0)$, $\Phi_\varrho^{-1}$ commutes with $\tilde\lambda$-shifts, $\Psi^{(n)}$ is invariant under $2\pi\Z^n$-shifts of $\lambda$ with injectivity modulo them (Theorem~\ref{bpP:p7:product-chart}(a)); $\mathcal T_1$ and $\Phi_\varrho$ are injective. This is Theorem~\ref{bpP:p9:diagonalize}(d) with $\mathcal T_1$.
\proofstep{(d) Mean}\label{stepP:p11:hc-mean} The $\theta$-average of $F$ is, by \Nref{bpN:p4:coefficients}(a) (the zero Fourier coefficient is holomorphic in the remaining variables) and Theorem~\ref{bpP:p9:diagonalize}(b), $S_{\rm fix}(\Lambda,\varrho x,\varrho y)=C_0+\sum_j\Omega_j\tilde R_j+O(\abs{(x,y)}^4)$ with $\tilde R_j=\frac12(x_j^2+y_j^2)$; the bound with explicit constant follows from Cauchy on $B_{r_z}$ after enlarging $M_F$: the function $R_4=F_0-C-\sum\Omega_j\tilde R_j$ is holomorphic with vanishing jet to order $3$ at $z=0$ (Corollary~\ref{bpP:p8:evenness}: even in $(\xi,\eta)$, no linear term, quadratic part as stated), hence $\abs{R_4}\le\sup\abs{R_4}(\abs z_\infty/r_z)^4$ (Schwarz lemma on polydiscs applied on the lines $z=tz_0$).
\end{proof}

\begin{lemma}[Total chart and lattice compatibility]\label{bpP:p11:total-chart}
Let $\Theta^{(1)}(\zeta,\tilde\rho)=\sqrt{\tilde\rho/\pi}\,(\sin2\pi\zeta,\cos2\pi\zeta)$ componentwise for $\tilde\rho_j>0$ (so $x_j=\sqrt{\tilde\rho_j/\pi}\sin2\pi\zeta_j$, $y_j=\sqrt{\tilde\rho_j/\pi}\cos2\pi\zeta_j$, $\frac12(x_j^2+y_j^2)=\tilde\rho_j/2\pi$). Then $\hat\Theta(\theta,J,\zeta,\tilde\rho)=(\theta,J,\Theta^{(1)}(\zeta,\tilde\rho))$ is a symplectic $1$-chart for $dJ\wedge d\theta+d\tilde\rho\wedge d\zeta$ (the $\mathcal T$-conjugate of the Part~7 polar chart $\Theta_n$), real analytic on $\{\tilde\rho_j>0\}$, $\hat\Theta(\theta,J,\zeta+k,\tilde\rho)=\hat\Theta(\theta,J,\zeta,\tilde\rho)$, and if $\hat\Theta(\cdot)=\hat\Theta(\cdot')$ then $\zeta'-\zeta\in\Z^n$ and the other coordinates agree. Consequently, for every open $\mathcal O\subset\{(\theta,J,\zeta,\tilde\rho):\tilde\rho_j>0\}$ with $\hat\Theta(\mathcal O)\subset\mathcal O_c$ and $\mathcal O+\iota_k=\mathcal O$ ($\iota_k=((k,0),0)$ shifting $(\theta,\zeta)$ by $k\in\Z^{2n}$), the chart $\Xi_0:=\Psi_c\circ\hat\Theta$ satisfies (Q1)--(Q2) of Theorem~\ref{bpP:p1:torus-transport} for $d=2n$ with $\sigma_k=0$.
\end{lemma}
\begin{proof}
\proofstep{Polar chart}\label{stepP:p11:tc-polar} Lemma~\ref{bpP:p7:lattice}(a),(b),(c) (with $\psi=2\pi\zeta$, $R=\tilde\rho/2\pi$) and Remark~\ref{bpP:p1:typical}(1),(3).
\proofstep{Lattice}\label{stepP:p11:tc-lat} $\Xi_0(w+\iota_k)=\Psi_c(\hat\Theta w)$ since $\hat\Theta$ is $\Z^n$-periodic in $\zeta$ and $\Psi_c(\cdot+(k_\theta))=\Psi_c(\cdot)$ (Proposition~\ref{bpP:p11:holo-chart}(c)), (Q1) with $\sigma_k=0$. (Q2): $\Xi_0(w')=\Xi_0(w)$ gives $\hat\Theta w'=\hat\Theta w+\iota_{(m,0)}$ by (c) of the proposition, then $\zeta'-\zeta\in\Z^n$ by the polar statement, $\tilde\rho'=\tilde\rho$, $J'=J$, so $w'=w+\iota_m$.
\end{proof}

\subsection{Step C: fast averaging (CP09 Lemma 4)}\label{bpP:p11:sec-C}

Fix $\nu_1:=3$. For $0<\varepsilon<1$ with $\varepsilon M_F<1$ put, with $J_0,\gamma_0,\tau_0,r_C,\sigma_F,r_z,M_F,L$ of Proposition~\ref{bpP:p11:holo-chart},
\begin{equation}\label{bpP:p11:eq:par1}
K_1=\frac{3(\nu_1-1)}{\pi\sigma_F}\log\frac1{\varepsilon M_F},\quad\alpha_1=\tfrac12\gamma_0K_1^{-\tau_0},\quad s_1=\min\Bigl\{r_C,\ \frac{\gamma_0}{4L\,K_1^{\tau_0+1}}\Bigr\},\quad d_1=\min\{s_1\sigma_F,\ r_z^2\}.
\end{equation}
(So $e^{-2\pi K_1\sigma_F/6}=(\varepsilon M_F)^{\nu_1-1}$, $\kappa_1:=2\pi K_1\sigma_F\ge6$ for small $\varepsilon$, and $s_1\asymp(\log\frac1\varepsilon)^{-(\tau_0+1)}$, $\alpha_1\asymp(\log\frac1\varepsilon)^{-\tau_0}$, $d_1\asymp(\log\frac1\varepsilon)^{-(\tau_0+1)}$.) Let $\mathcal D_1=\T_{\sigma_F}\times D(J_0,s_1)\times B_{r_z}$ and $\mathcal D_1'=\T_{\sigma_F/6}\times D(J_0,s_1/2)\times B_{r_z/2}$.

\begin{proposition}[Fast averaging with the secular pairs as spectators]\label{bpP:p11:avg-fast}
There is $\varepsilon_1\in(0,1)$ depending only on the data such that for $0<\varepsilon\le\varepsilon_1$ the inequality $(S1)$: $2^{12}(1+n)K_1\,\varepsilon M_F\,\sigma_F\le\alpha_1d_1$ holds, and there is a holomorphic, injective, real, symplectic map $\Phi^1_\varepsilon:\mathcal D_1'\to\mathcal D_1$ for $dJ\wedge d\theta+dy\wedge dx$ ($\Phi^1_\varepsilon-\mathrm{id}$ $\Z^n$-periodic in $\theta$) with
\[H\circ\Phi^1_\varepsilon=\tilde h(J)+\varepsilon G(J,z)+f_1(\theta,J,z)\quad\text{on }\mathcal D_1',\qquad H:=\tilde h+\varepsilon F,\quad G:=F_0+\varepsilon\bar g,\]
where $G=G(J,z)$ is holomorphic and real, $f_1$ is real, and
\[\sup_{D(J_0,s_1/2)\times B_{r_z/2}}\abs{\varepsilon\bar g}\le\eta_1(\varepsilon):=\frac{400(1+n)\,\varepsilon M_F^2}{2\pi\alpha_1d_1},\qquad\wn{f_1}{D(J_0,s_1/2)\times B_{r_z/2}}{\sigma_F/6}\le(\varepsilon M_F)^{\nu_1},\]
and the displacement bounds $2\pi\alpha_1\sigma_F\abs{J-J'}\le3\varepsilon M_F$, $2\pi\alpha_1s_1\abs{\theta-\theta'}\le9\varepsilon M_F$, $2\pi\alpha_1r_z\abs{z-z'}\le9\varepsilon M_F$ for $(\theta',J',z')\in\mathcal D_1'$, $(\theta,J,z)=\Phi^1_\varepsilon(\theta',J',z')$. Note $\eta_1(\varepsilon)\le c\,\varepsilon(\log\frac1\varepsilon)^{2\tau_0+1}$ with $c=c(\mathrm{data})$ (the exponent $2\tau_0+1$ is that of CP09 (25), $\Omega-\hat\Omega=O(\varepsilon\log^{2\tau_0+1}\varepsilon^{-1})$).
\end{proposition}
\begin{proof}
\proofstep{Application of Part 2}\label{stepP:p11:af-apply} Proposition~\ref{bpP:p2:osc-averaging} with $n\to n$ (angles), $m\to n$ (oscillator pairs; Part~2's form is $d\theta\wedge dI+dq\wedge dp$ with $(q,p)=(x,y)$ the position and momentum of the pair, which is the \emph{negative} of $dJ\wedge d\theta+dy\wedge dx$; a map preserves one if and only if it preserves the other, and the Hamilton equations $\dot\theta=\partial_JH$, $\dot J=-\partial_\theta H$, $\dot x=\partial_yH$, $\dot y=-\partial_xH$ of $X_u$ in Part~2 are the frozen ones), $D=\{J_0\}$, $r=s_1$, $r_p=r_q=r_z$, $\bar s=0$, $s=\sigma_F$, $K=K_1$, $\alpha=\alpha_1$, $h=\tilde h$ (independent of $z$), $f=\varepsilon F$, $E=\varepsilon M_F$ (Proposition~\ref{bpP:p11:holo-chart}(a)), $b=1+n$, $d=d_1$; $\kappa=\kappa_1\ge6$.
\proofstep{(O1)}\label{stepP:p11:af-o1} For $J\in D(J_0,s_1)$ and $0<\abs k_1\le K_1$: $\abs{k\cdot\omega(J)}\ge\abs{k\cdot\omega(J_0)}-\abs k_1\abs{\omega(J)-\omega(J_0)}_\infty\ge\gamma_0K_1^{-\tau_0}-K_1Ls_1\ge\tfrac34\gamma_0K_1^{-\tau_0}\ge\alpha_1$ (Lemma~\ref{bpP:p11:select}, \Kref{bp:p1:ext-mvt}, definition of $s_1$).
\proofstep{(O2) is (S1)}\label{stepP:p11:af-o2} $K_1,\alpha_1^{-1},d_1^{-1}$ are polynomials in $\log\frac1\varepsilon$, so the ratio $K_1\varepsilon M_F\sigma_F/(\alpha_1d_1)\to0$ (e4); $\varepsilon_1$ is the largest value with (S1) and $\kappa_1\ge6$ for all smaller $\varepsilon$.
\proofstep{Conclusion}\label{stepP:p11:af-concl} Proposition~\ref{bpP:p2:osc-averaging} gives $\Phi^1_\varepsilon$ and $H\circ\Phi^1_\varepsilon=\tilde h+g+f_*$, $g=g(J,z)$ independent of $\theta$, with $\lVert g-\hat f_0\rVert\le400bE^2/(\bar\alpha d)$, $\bar\alpha=2\pi\alpha_1$, $\hat f_0=\varepsilon F_0$; put $\varepsilon^2\bar g=g-\varepsilon F_0$ and $f_1=f_*$, $\lVert f_*\rVert\le e^{-\kappa_1/6}E=(\varepsilon M_F)^{\nu_1-1}\varepsilon M_F$; the displacement bounds are those of that proposition (with $\bar\alpha s_1\abs{\theta-\theta'}$ replaced by $\bar\alpha r\abs{\theta-\theta'}$, $r=s_1$). The sup bound for $\varepsilon\bar g$ is \Nref{bpN:p4:sup-weighted}(a).
\end{proof}

\subsection{Step D: polar normalisation and slow averaging}\label{bpP:p11:sec-D}

CP09 now passes to normal form coordinates and normalises the quadratic and higher Taylor terms of $G$ in $z$ up to a finite order. Here the dependence on $z$ is kept and removed in action--angle variables on the shell $\tilde\rho_j\approx2\pi\varepsilon$, by averaging over the polar angles. That needs one extension of Part~2's two-scale averaging.

\begin{lemma}[Blockwise bracket and field estimates]\label{bpP:p11:aniso-bracket}
Let $n=n_1+n_2$, radii $r^0_{(j)}\ge r_{(j)}>\rho_{(j)}>0$ ($j=1,2$) and $V^0\supset V\supset V-\rho$ the corresponding products of polydiscs in $\C^n$ (block $j$ has radius $r^0_{(j)}$, $r_{(j)}$, $r_{(j)}-\rho_{(j)}$), $\rho_*=\min_j\rho_{(j)}$, $0<\sigma<s\le s_0$. For $u$ holomorphic on $\T_{s_0}\times V^0$ and $w$ on $\T_s\times V$ with finite weighted norms, $\{u,w\}=\sum_i(\partial_{\theta_i}u\,\partial_{I_i}w-\partial_{I_i}u\,\partial_{\theta_i}w)$ satisfies
\[\lVert\{u,w\}\rVert_{V-\rho,s-\sigma}\le\frac{\nu_\theta+\nu_I}{e\,\rho_*\sigma}\lVert u\rVert_{V^0,s_0}\lVert w\rVert_{V,s},\]\[\nu_\theta=\frac\sigma{s_0-s+\sigma},\qquad\nu_I=\frac{\rho_*}{\min_j(r^0_{(j)}-r_{(j)}+\rho_{(j)})}\le1,\]
and for $\phi$ on $\T_s\times V$: on $\T_{s-\sigma}\times(V-\rho)$ the vector field $X_\phi=(\partial_I\phi,-\partial_\theta\phi)$ satisfies $\nu(X_\phi)\le\lVert\phi\rVert_{V,s}/(\rho_*\sigma)$ for the scaled norm $\nu(a,b)=\max\{\lvert a\rvert_\infty/\sigma,\ \lvert b_{(j)}\rvert_\infty/\rho_{(j)}\}$ (whose unit balls are the polydiscs of the losses).
\end{lemma}
\begin{proof}
\proofstep{Coordinate derivatives}\label{stepP:p11:ab-der} Fix $I\in V-\rho$ and apply \Nref{bpN:p4:cauchy}(a),(b) to each variable separately: $\sum_i\lVert\partial_{\theta_i}u\rVert_{s-\sigma}\le\lVert u\rVert_{s_0}/(e(s_0-s+\sigma))$ at fixed $I$ (the loss is to angular width $s-\sigma$; all bracket factors of Step~\ref{stepP:p11:ab-prod} are evaluated at $I\in V-\rho$ and angular width $s-\sigma$) (angular Cauchy at fixed parameter, \Nref{bpN:p4:cauchy}(a)); for the action derivative of $w$ in a coordinate of block $j$, the Cauchy disc has radius $\ge\rho_{(j)}\ge\rho_*$, so $\max_i\lvert\partial_{I_i}w\rvert\le\lVert w\rVert/\rho_*$ (global norm, \Nref{bpN:p4:cauchy}(b)); for $u$ the disc radius is $\ge r^0_{(j)}-r_{(j)}+\rho_{(j)}$.
\proofstep{Products}\label{stepP:p11:ab-prod} By the algebra property \Nref{bpN:p4:algebra} at each fixed $I$ and the triangle inequality, $\lVert\{u,w\}\rVert\le\bigl(\sum_i\lVert\partial_{\theta_i}u\rVert\bigr)\max_i\lVert\partial_{I_i}w\rVert+\max_i\lVert\partial_{I_i}u\rVert\bigl(\sum_i\lVert\partial_{\theta_i}w\rVert\bigr)\le\frac{\lVert u\rVert\lVert w\rVert}e\Bigl(\frac1{(s_0-s+\sigma)\rho_*}+\frac1{\min_j(r^0_{(j)}-r_{(j)}+\rho_{(j)})\,\sigma}\Bigr)$, which is the claim. The field bound: $\lvert\partial_{I_i}\phi\rvert/\sigma\le\lVert\phi\rVert/(\rho_{(j)}\sigma)$ and $\lvert\partial_{\theta_i}\phi\rvert/\rho_{(j)}\le\lVert\phi\rVert/(e\sigma\rho_{(j)})$, both $\le\lVert\phi\rVert/(\rho_*\sigma)$.
\end{proof}

\begin{lemma}[Two-scale averaging with two action radii]\label{bpP:p11:aniso-avg}
Let $n=n_1+n_2$, $I^0\in\R^n$, $r_1,r_2>0$, $r_*=\min\{r_1,r_2\}$, $V=\{I\in\C^n:\abs{I_i-I_i^0}<r_1\ (i\le n_1),\ \abs{I_i-I^0_i}<r_2\ (i>n_1)\}$, $s>0$, $K>0$ with $\kappa=2\pi Ks\ge6$, $0<\alpha_2\le\alpha_1$, $\bar\alpha_j=2\pi\alpha_j$. Let $h$ be holomorphic on $V$ and real on real points, $\omega=\partial_Ih$, with
\begin{enumerate}[label=(B\arabic*)]
\item $\abs{k\cdot\omega(I)}\ge\alpha_1$ for $k_1\ne0$ and $\ge\alpha_2$ for $k_1=0$, $k\ne0$, $I\in V$, $\abs k_1\le K$ ($k=(k_1,k_2)\in\Z^{n_1}\times\Z^{n_2}$);
\item $f$ holomorphic on $\T_{s}\times V$, $1$-periodic, real on real points, $E:=\wn fVs\le\alpha_2r_*/(128K)$.
\end{enumerate}
Then the conclusions (a)--(d) of Proposition~\ref{bpP:p2:nfb} hold with $D_{r}$ replaced by $V$, $\Dom{r/2}{\bar s+s/6}$ by $\T_{s/6}\times\tfrac12V$ ($\tfrac12V$: radii $r_1/2,r_2/2$), $\bar s=0$, with $\bar\alpha_2r$ in (d) replaced by $\bar\alpha_2r_*$, and in addition
\[\wn{g-\hat f_0}{\frac12V}{s/6}\le\frac{130\,E^2}{\bar\alpha_2r_*s},\qquad g=g(I')\ \text{real},\quad H\circ\Psi=h+g+f_*,\quad\wn{f_*}{\frac12V}{s/6}\le e^{-\kappa/6}E.\]
\end{lemma}
\begin{proof}
\proofstep{The proof of Part 2, with the named replacements}\label{stepP:p11:aa-same} Proposition~\ref{bpP:p2:nfb} is proved from Lemma~\ref{bpP:p2:direct}, Lemma~\ref{bpP:p2:split} and \Nref{bpN:p4:bracket}, \Nref{bpN:p5:vf-bound}, \Nref{bpN:p5:flow}, \Nref{bpN:p5:composition}, \Nref{bpN:p5:transform}, \Nref{bpN:p5:canonical-map}, \Nref{bpN:p6:homological}, \Nref{bpN:p6:iterative}, \Nref{bpN:p6:telescope}. The action radii enter these only through (i) the bracket and field bounds and (ii) the weights $\rho/(r_j-r+\rho)$ of $\lvert g\rvert_*$ (\Nref{bpN:p6:setting}); the divisor hypothesis (B1) is a statement on all of $V$ and is block independent, and the homological solution $\phi$ is a function on $\T_s\times V$ with $\lVert\phi\rVert\le E/\bar\alpha_2$, $\lVert\phi^{(1)}\rVert\le E/\bar\alpha_1$ exactly as in Lemma~\ref{bpP:p2:split} (which uses \Nref{bpN:p5:vf-bound} only through $\lvert\partial_\theta\phi\rvert$, $\lvert\partial_I\phi\rvert$ and now Lemma~\ref{bpP:p11:aniso-bracket}). The replacements: (R1) $\Nref{bpN:p4:bracket}$ by Lemma~\ref{bpP:p11:aniso-bracket} with $\rho\sigma\to\rho_*\sigma$ (the weight $\nu_I\le1$ replaces $\rho/(r_0-r+\rho)$); (R2) \Nref{bpN:p5:vf-bound}, \Nref{bpN:p5:flow}, \Nref{bpN:p5:composition}, \Nref{bpN:p5:transform}, \Nref{bpN:p5:canonical-map} by the same statements for the scaled norm $\nu$ of Lemma~\ref{bpP:p11:aniso-bracket} (these results are stated for an arbitrary scaled norm and use the geometry only through the two estimates of that lemma, Lemma~\ref{bpP:p2:osc-iter}, step~\ref{stepP:p2:oi-transfer}); (R3) every product $\bar\alpha\rho\sigma$ by $\bar\alpha\rho_*\sigma$, with the losses chosen block by block: step $0$ has $\rho_{(j),0}=r_{(j)}/8$, $\sigma_0=s/6$, hence $y_0=E/(\bar\alpha_2\rho_{*,0}\sigma_0)=48E/(\bar\alpha_2r_*s)$ ($\rho_{*,0}=r_*/8$); steps $i\ge1$ have $\rho_{(j)}=r_{(j)}/(8N')$ ($N'$ the number of steps, $N'=\lfloor\kappa/(12\ln2)\rfloor$) and radii $\frac34r_{(j)}-2(i-1)\rho_{(j)}$, so that all ratios $\rho_{(j)}/(r_{(j)}-r_{(j),i}+\rho_{(j)})$ are identical to those of the isotropic proof and $\rho_*=r_*/(8N')$. The inductive inequalities (2.2) and the claim $\varepsilon_{i+1}\le\varepsilon_i/4$ of Step~\ref{stepP:p2:nfb-ind} are numerical consequences of $y_0$ and of $\varepsilon\le\bar\alpha_2\rho_*\sigma\cdot(\ldots)$, i.e.\ of (B2) with $r_*$ in place of $r$, and are unchanged. The displacements of $I_j$ ($\le E/(e\bar\alpha_j\sigma)$) do not involve any radius; that of $\theta$ is $\le E/(\bar\alpha_2\rho_*)$ per step, giving $9E/(\bar\alpha_2r_*)$ as in Step~\ref{stepP:p2:nfb-disp} with $r\to r_*$.
\proofstep{The quadratic bound for $g$}\label{stepP:p11:aa-quad} In Step~\ref{stepP:p2:nfb-first} (Case A) one has $\varepsilon_1:=\lVert f_1\rVert\le\frac87(1+\frac2e)y_0E$, $y_0=48E/(\bar\alpha_2r_*s)$, so $\varepsilon_1\le1.984\cdot48E^2/(\bar\alpha_2r_*s)\le96E^2/(\bar\alpha_2r_*s)$; Step~\ref{stepP:p2:nfb-concl} gives $\lVert g-g_0\rVert\le\frac43\varepsilon_1\le128E^2/(\bar\alpha_2r_*s)$. In Case B, $g=\hat f_0$. (Part~2 records only the consequence $\le E/\kappa$ of this bound.)
\end{proof}
\sourcegap{Lemma~\ref{bpP:p11:aniso-avg} is a variant of Proposition~\ref{bpP:p2:nfb} (two action radii), obtained by the same list of changes as Lemma~\ref{bpP:p2:osc-iter}; it is the statement that Part~11 needs and is not in CP10/CP09 (CP09 Lemma 4 cites [3, Prop.~A.1] with the single radius $s$). It may be absorbed into Part~2.}

For $0<\varepsilon\le\varepsilon_1$ let $\Pi(\zeta,\tilde\rho):=\Theta^{(1)}(\zeta,\tilde\rho)$ (Lemma~\ref{bpP:p11:total-chart}) on the \emph{polar shell}
\[\mathcal P_\varepsilon=\{(\zeta,\tilde\rho)\in\C^n\times\C^n:\ \abs{\operatorname{Im}\zeta}_\infty<\sigma_F/6,\ \abs{\tilde\rho_j-2\pi\varepsilon}<\pi\varepsilon\ (\forall j)\}.\]
$\Pi$ is holomorphic on $\mathcal P_\varepsilon$ ($\operatorname{Re}\tilde\rho_j>\pi\varepsilon>0$, principal square root) and $\abs{\Pi(\zeta,\tilde\rho)}_\infty\le\sqrt{3\varepsilon}\,e^{\pi\sigma_F/3}$; hence $\Pi(\mathcal P_\varepsilon)\subset B_{r_z/4}$ as soon as $\sqrt{3\varepsilon}e^{\pi\sigma_F/3}\le r_z/4$, which we assume ($\varepsilon\le\varepsilon_2$). Put $\sigma_3=\sigma_F/12$ and, for $\varepsilon\le\varepsilon_2$,
\begin{equation}\label{bpP:p11:eq:par2}
\begin{aligned}
&\hat G(J,\zeta,\tilde\rho)=G(J,\Pi(\zeta,\tilde\rho)),\quad\hat G_0=\textstyle\int_{\T^n}\hat G\,d\zeta,\quad h_0(J,\tilde\rho)=\tilde h(J)+\varepsilon\hat G_0(J,\tilde\rho),\\ &f_2=\varepsilon(\hat G-\hat G_0)+f_1\circ(\mathrm{id},\Pi),\qquad K_2=\frac{3\nu_2}{\pi\sigma_3}\log\frac1\varepsilon,\quad\alpha_{(1)}=\tfrac14\gamma_0K_2^{-\tau_0},\quad\alpha_{(2)}=\tfrac14\varepsilon\gamma_0K_2^{-\tau_0},\\ &s_2=\min\Bigl\{\frac{s_1}4,\ \frac{\gamma_0}{8L\,K_2^{\tau_0+1}}\Bigr\},\quad r_*=\pi\varepsilon .\\
\end{aligned}
\end{equation}
Here $\nu_2\ge1$ is a real number fixed in Theorem~\ref{bpP:p11:wrapper}; $e^{-2\pi K_2\sigma_3/6}=\varepsilon^{\nu_2}$. The slow frequency of $\tilde\Omega$ is the second block: $\partial_{\tilde\rho}h_0=\varepsilon(\tilde\Omega(J)+\cdots)$ with $\tfrac12(x_j^2+y_j^2)=\tilde\rho_j/2\pi$ and $\Omega_j/2\pi=\tilde\Omega_j$ (Proposition~\ref{bpP:p11:holo-chart}(d)).

\begin{lemma}[Polar average]\label{bpP:p11:polar-avg}
Let $0<r\le1$, $S>0$, $D\subset\C^n$ open and $G:D\times B_r\to\C$ holomorphic, $B_r=D(0,r)\subset\C^{2n}$ ($z=(x,y)$), $\lvert G\rvert\le S$. Then there is a holomorphic $\mathcal AG:D\times D(0,\pi r^2/4)\to\C$ (depending on $G$ only) such that
\begin{enumerate}[label=(\alph*)]
\item $\lvert\mathcal AG\rvert\le2^{2n}S$ on $D\times D(0,\pi r^2/4)$ (this is the radius $r'=r/2$ below; one has the same lemma with $r'$ arbitrary in $(0,r)$, constant $(1-r'/r)^{-2n}$ and radius $\pi r'^2$);
\item for $J\in D$, real $\zeta\in\R^n$ and $\tilde\rho$ with $\operatorname{Re}\tilde\rho_j>0$, $\lvert\tilde\rho_j\rvert<\pi r^2/4$: $\int_{\T^n}G(J,\Pi(\zeta,\tilde\rho))\,d\zeta=\mathcal AG(J,\tilde\rho)$ (principal square roots);
\item if the Taylor coefficients of $G(J,\cdot)$ of degree $<k$ vanish identically in $J$, then $\lvert\mathcal AG(J,\tilde\rho)\rvert\le2^{2n}S\bigl(4\lvert\tilde\rho\rvert_\infty/(\pi r^2)\bigr)^{\lceil k/2\rceil}$;
\item $\mathcal A$ is linear, $\mathcal A(C)=C$ for $z$-independent $C$, and $\mathcal A\bigl(\tfrac12(x_j^2+y_j^2)\bigr)=\tilde\rho_j/2\pi$.
\end{enumerate}
\end{lemma}
\begin{proof}
\proofstep{Taylor series}\label{stepP:p11:pa-taylor} Write $G(J,z)=\sum_{a,b\in\N^n}c_{ab}(J)x^ay^b$. By (e2) on $B_r$, $\lvert c_{ab}(J)\rvert\le Sr^{-\lvert a\rvert-\lvert b\rvert}$, so $\sum_{a,b}\lvert c_{ab}\rvert(r/2)^{\lvert a\rvert+\lvert b\rvert}\le2^{2n}S$ (e2).
\proofstep{Averaging a monomial}\label{stepP:p11:pa-mono} For real $\zeta$ and $\lvert\tilde\rho_j\rvert<\pi r^2/4$ (so $\lvert\Pi_j\rvert\le\lvert\tilde\rho_j/\pi\rvert^{1/2}<r/2$ componentwise, as $\lvert\sin\rvert,\lvert\cos\rvert\le1$ on reals): $x^ay^b\circ\Pi=\prod_j(\tilde\rho_j/\pi)^{(a_j+b_j)/2}\sin^{a_j}(2\pi\zeta_j)\cos^{b_j}(2\pi\zeta_j)$ and by (e5) its $\zeta$-integral is $0$ unless all $a_j,b_j$ are even, in which case it is $\prod_jm_j(\tilde\rho_j/\pi)^{\kappa_j}$ with $\kappa_j=(a_j+b_j)/2\in\N$, $\lvert m_j\rvert\le1$; thus a polynomial in $\tilde\rho$ (no square roots remain: this is the cancellation that the polar inverse, singular at $\tilde\rho_j=0$, would otherwise obstruct).
\proofstep{The series}\label{stepP:p11:pa-series} Put $\mathcal AG(J,\tilde\rho):=\sum_{a,b\ \mathrm{even}}c_{ab}(J)\,m_{ab}\,\pi^{-\lvert\kappa\rvert}\tilde\rho^{\kappa}$, $m_{ab}=\prod_jm_j(a_j,b_j)$. For $\lvert\tilde\rho_j\rvert<\pi r^2/4$ one has $\lvert\tilde\rho^\kappa\pi^{-\lvert\kappa\rvert}\rvert\le(r/2)^{\lvert a\rvert+\lvert b\rvert}$, so the series converges normally by Step~\ref{stepP:p11:pa-taylor}, defines a holomorphic function of $(J,\tilde\rho)$ ($c_{ab}$ holomorphic in $J$, uniform limit of holomorphic partial sums), with the bound (a). (b): termwise integration of the uniformly convergent series (e5). (c): only $\lvert a\rvert+\lvert b\rvert\ge k$, hence $\lvert\kappa\rvert\ge\lceil k/2\rceil$, occurs; bound each term by $\lvert c_{ab}\rvert(r/2)^{\lvert a\rvert+\lvert b\rvert}\bigl(4\lvert\tilde\rho\rvert_\infty/(\pi r^2)\bigr)^{\lvert\kappa\rvert}$ and use $4\lvert\tilde\rho\rvert_\infty/(\pi r^2)<1$. (d): $\frac12(x_j^2+y_j^2)\circ\Pi=\tilde\rho_j/2\pi$ is already $\zeta$-independent.
\end{proof}

\begin{proposition}[Polar normalisation and slow averaging]\label{bpP:p11:avg-slow}
Let $\nu_2\ge1$. There are $\varepsilon_2\in(0,\varepsilon_1]$ and $c_2\ge1$ depending only on the data and $\nu_2$ such that for $0<\varepsilon\le\varepsilon_2$, with $V_2=D(J_0,s_2)\times D_{\tilde\rho}$ ($D_{\tilde\rho}=\{\abs{\tilde\rho_j-2\pi\varepsilon}<\pi\varepsilon\}$, radii $s_2$ for $J$ and $r_2=\pi\varepsilon$ for $\tilde\rho$), $n_1=n_2=n$, angles $(\theta,\zeta)$, $s=\sigma_3$:
\begin{enumerate}[label=(\alph*)]
\item (\emph{the perturbation}) $\wn{f_2}{V_2}{\sigma_3}\le E_2:=c_2\,\varepsilon^{5/2}(\log\frac1\varepsilon)^{2\tau_0+1}$, and $\hat G_0=C(J)+\sum_j\tilde\Omega_j(J)\tilde\rho_j+\hat Z(J,\tilde\rho)$ with $\hat Z$ holomorphic on $D(J_0,s_1/2)\times D(0,\pi r_z^2/16)$ and $\abs{\partial_J^\alpha\partial_{\tilde\rho}\hat Z}\le c_2(\varepsilon+\eta_1(\varepsilon))\alpha!\,(4/s_1)^{\abs\alpha}$ on $D(J_0,s_1/4)\times D_{\tilde\rho}$;
\item (\emph{(B1)--(B2) for $h_0$ and $f_2$}) (B1) holds on $V_2$ with $\alpha_1\to\alpha_{(1)}$, $\alpha_2\to\alpha_{(2)}$ and $K=K_2$, and (B2) holds: $E_2\le\alpha_{(2)}\pi\varepsilon/(128K_2)$;
\item (\emph{averaging}) Lemma~\ref{bpP:p11:aniso-avg} applies to $(h_0,f_2)$ and gives $\Psi_2:\T_{\sigma_3/6}\times\tfrac12V_2\to\T_{\sigma_3}\times V_2$ with $(H\circ\Phi^1_\varepsilon\circ\hat\Theta)\circ\Psi_2=h_0+g_2+f_3$, where $\hat\Theta=(\mathrm{id}_{\theta,J},\Pi)$, $g_2=g_2(J,\tilde\rho)$ is real and
\[\abs{g_2}\le c_2\,\varepsilon^{3}(\log\tfrac1\varepsilon)^{5\tau_0+2}\ \text{ on }\tfrac12V_2,\qquad\wn{f_3}{\frac12V_2}{\sigma_3/6}\le\varepsilon^{\nu_2}E_2;\]
\item (\emph{displacements}) for $(\theta,\zeta,J,\tilde\rho)=\Psi_2(\theta',\zeta',J',\tilde\rho')$: $\abs{\tilde\rho-\tilde\rho'}\le c_2\varepsilon^{3/2}(\log\frac1\varepsilon)^{3\tau_0+1}$, $\abs{J-J'}\le c_2\varepsilon^{5/2}(\log\frac1\varepsilon)^{3\tau_0+1}$, $\abs{\theta-\theta'},\abs{\zeta-\zeta'}\le c_2\varepsilon^{1/2}(\log\frac1\varepsilon)^{3\tau_0+1}$.
\end{enumerate}
\end{proposition}
\begin{proof}
\proofstep{(a) Size of $\hat G-\hat G_0$ and $\hat Z$}\label{stepP:p11:as-osc} $G=C+\sum_j\Omega_j\tilde R_j+R_4+\varepsilon\bar g$ with $\tilde R_j=\tilde\rho_j/2\pi$ independent of $\zeta$, on $D(J_0,s_1/2)\times B_{r_z/2}$. Apply Lemma~\ref{bpP:p11:polar-avg} with $r=r_z/2$ to $R_4+\varepsilon\bar g$: $\hat Z:=\mathcal A(R_4+\varepsilon\bar g)$ is holomorphic on $D(J_0,s_1/2)\times D(0,\pi r_z^2/16)$ and $\hat G_0=C+\sum_j\tilde\Omega_j(J)\tilde\rho_j+\hat Z$ for $\operatorname{Re}\tilde\rho_j>0$, $\lvert\tilde\rho_j\rvert<\pi r_z^2/16$ (Lemma~\ref{bpP:p11:polar-avg}(b),(d) and $\Omega_j/2\pi=\tilde\Omega_j$), in particular on $D_{\tilde\rho}$ for $\varepsilon\le\varepsilon_2$. \emph{Derivatives.} $R_4$ has vanishing jet to order $3$ and $\lvert R_4\rvert\le M_F/16$ on $B_{r_z/2}$ (Proposition~\ref{bpP:p11:holo-chart}(d)), so by Lemma~\ref{bpP:p11:polar-avg}(a),(c) $\lvert\mathcal AR_4\rvert\le2^{2n}M_F(16\lvert\tilde\rho\rvert_\infty/(\pi r_z^2))^2/16$; at points of $D_{\tilde\rho}$ ($\lvert\tilde\rho\rvert_\infty\le3\pi\varepsilon$) the Cauchy estimate (e2) in $\tilde\rho$ on the polydisc of radius $\pi\varepsilon/4$ gives $\lvert\partial_{\tilde\rho}\mathcal AR_4\rvert\le c\,M_F\,\varepsilon\,r_z^{-4}$; $J$-derivatives of order $\alpha$ cost $\alpha!(4/s_1)^{\lvert\alpha\rvert}$ (Cauchy on radius $s_1/4$ inside $D(J_0,s_1/2)$). For $\varepsilon\bar g$ (bounded by $\eta_1$, Proposition~\ref{bpP:p11:avg-fast}) $\lvert\mathcal A(\varepsilon\bar g)\rvert\le2^{2n}\eta_1$, so by (e2) with radius $\pi r_z^2/32$: $\lvert\partial_{\tilde\rho}\mathcal A(\varepsilon\bar g)\rvert\le c\,\eta_1r_z^{-2}$, with the same $J$-losses. This is the bound of (a) for $\hat Z$. \emph{Oscillation.} $\hat G-\hat G_0=(R_4+\varepsilon\bar g)\circ\Pi-\mathcal A(R_4+\varepsilon\bar g)$ on $\mathcal P_\varepsilon\times D(J_0,s_1/2)$ (the identity (b) of the lemma holds for real $\zeta$ and extends to complex $\zeta$ because $\hat G_0$ does not depend on $\zeta$ and the $\zeta$-Fourier mean of a function holomorphic in $\zeta$ is the same). On $\mathcal P_\varepsilon$: $\lvert\Pi\rvert_\infty\le\sqrt{3\varepsilon}e^{\pi\sigma_F/3}\le r_z/4$, so $\lvert R_4\circ\Pi\rvert\le M_F(\lvert\Pi\rvert/r_z)^4\le9e^{4\pi\sigma_F/3}M_F\varepsilon^2r_z^{-4}$; and on the inner polydisc $\lvert z\rvert_\infty\le r_z/4$ the Cauchy estimate for $\varepsilon\bar g$ (bounded by $\eta_1$ on $B_{r_z/2}$, radius $r_z/4$ around each point) gives $\lvert\partial_{z_i}(\varepsilon\bar g)\rvert\le4\eta_1/r_z$, so $\lvert\varepsilon\bar g(J,z)-\varepsilon\bar g(J,0)\rvert\le2n\cdot4\eta_1\lvert z\rvert_\infty/r_z$ (mean value, \Kref{bp:p1:ext-mvt}); also $\lvert\mathcal A(\varepsilon\bar g)-\varepsilon\bar g(J,0)\rvert\le\sup_\zeta\lvert\varepsilon\bar g\circ\Pi-\varepsilon\bar g(\cdot,0)\rvert$ (the mean of the difference). Hence the $\zeta$-oscillation of $\varepsilon\bar g\circ\Pi$ is $\le16n\eta_1\sqrt{3\varepsilon}e^{\pi\sigma_F/3}/r_z$ and $\sup\lvert\hat G-\hat G_0\rvert\le c(\varepsilon^2+\eta_1\sqrt\varepsilon)$ with $c$ depending on $n,M_F,r_z,\sigma_F$; times $\varepsilon$, the dominant part is $c\,\varepsilon^{5/2}(\log\frac1\varepsilon)^{2\tau_0+1}$ (the term $\varepsilon\cdot\varepsilon^2$ is smaller). The part $f_1\circ(\mathrm{id},\Pi)$ is bounded by $(\varepsilon M_F)^3$ on $\T_{\sigma_F/6}$ (Proposition~\ref{bpP:p11:avg-fast}; $\Pi(\mathcal P_\varepsilon)\subset B_{r_z/4}\subset B_{r_z/2}$). Passing from the supremum on the strip of width $\sigma_F/6=2\sigma_3$ to $\lVert\cdot\rVert_{V_2,\sigma_3}$ costs $\coth(\pi\sigma_3)^{2n}$ (\Nref{bpN:p4:sup-weighted}(b); both $\theta$ and $\zeta$ are angles, the oscillator coordinates are parameters). All constants are absorbed into $c_2$.
\proofstep{(b) Non-resonance on $V_2$}\label{stepP:p11:as-nr} $\partial_Jh_0=\omega(J)+\varepsilon\partial_J\hat G_0$ and $\partial_{\tilde\rho}h_0=\varepsilon(\tilde\Omega(J)+\partial_{\tilde\rho}\hat Z)$. By (e2) (Cauchy radius $s_1/4$, $V_2\subset D(J_0,s_1/4)\times D_{\tilde\rho}$ since $s_2\le s_1/4$) $\lvert\partial_J\hat G_0\rvert\le cM_F/s_1$, and by Step~\ref{stepP:p11:as-osc} $\lvert\partial_{\tilde\rho}\hat Z\rvert\le c_2(\varepsilon+\eta_1)$. Put $\delta_0=\gamma_0K_2^{-\tau_0}$. By Lemma~\ref{bpP:p11:select}, for $\lvert k_i\rvert_1\le K_2$: $\lvert k_1\cdot\omega(J_0)\rvert\ge\delta_0$ ($k_1\ne0$), $\lvert k_2\cdot\tilde\Omega(J_0)\rvert\ge\delta_0$ ($k_2\ne0$), and $\lvert k_i\cdot(\varpi_J(J)-\varpi_J(J_0))\rvert\le K_2Ls_2\le\frac18\delta_0$ for $J\in D(J_0,s_2)$. \emph{Pure slow case} $k_1=0$, $k_2\ne0$: $k\cdot\partial h_0=\varepsilon\bigl(k_2\cdot\tilde\Omega(J)+k_2\cdot\partial_{\tilde\rho}\hat Z\bigr)$, and \emph{without the prefactor $\varepsilon$} $K_2\lvert\partial_{\tilde\rho}\hat Z\rvert\le K_2c_2(\varepsilon+\eta_1)\le\frac18\delta_0$ for $\varepsilon\le\varepsilon_2$ (since $(\varepsilon+\eta_1)K_2^{\tau_0+1}\le c\,\varepsilon(\log\frac1\varepsilon)^{3\tau_0+2}\to0$, (e4)); hence $\lvert k\cdot\partial h_0\rvert\ge\varepsilon(\delta_0-\frac18\delta_0-\frac18\delta_0)\ge\alpha_{(2)}$. \emph{Case $k_1\ne0$}: $\lvert k\cdot\partial h_0\rvert\ge\lvert k_1\cdot\omega(J)\rvert-\varepsilon\lvert k_2\rvert_1\lvert\tilde\Omega\rvert-\varepsilon\lvert k\rvert_1\max\{\lvert\partial_J\hat G_0\rvert,\lvert\partial_{\tilde\rho}\hat Z\rvert\}\ge\frac78\delta_0-\varepsilon K_2(L+c(1+M_F/s_1))\ge\frac12\delta_0\ge\alpha_{(1)}$ for $\varepsilon\le\varepsilon_2$ (here the prefactor $\varepsilon$ is present and $\varepsilon K_2^{\tau_0+1}s_1^{-1}\to0$). $\alpha_{(2)}\le\alpha_{(1)}$.
\proofstep{(B2)}\label{stepP:p11:as-b2} $E_2\le c_2\varepsilon^{5/2}(\log\frac1\varepsilon)^{2\tau_0+1}$ and $\alpha_{(2)}\pi\varepsilon/(128K_2)=\frac{\pi\gamma_0}{512}\varepsilon^2K_2^{-\tau_0-1}$; the ratio of the former to the latter is $\varepsilon^{1/2}$ times a power of $\log\frac1\varepsilon$, so (B2) holds for $\varepsilon\le\varepsilon_2$ (e4). $\kappa=2\pi K_2\sigma_3=6\nu_2\log\frac1\varepsilon\ge6$.
\proofstep{(c),(d)}\label{stepP:p11:as-concl} Lemma~\ref{bpP:p11:aniso-avg} with $s=\sigma_3$, $r_*=\pi\varepsilon\le s_2$, $\bar\alpha_2=2\pi\alpha_{(2)}$: $g_2:=g$, $f_3:=f_*$, $e^{-\kappa/6}=\varepsilon^{\nu_2}$; $\hat f_0=\int f_2=\int_{\T^{2n}}f_1\circ(\mathrm{id},\Pi)$ (the term $\varepsilon(\hat G-\hat G_0)$ has zero $\zeta$-mean and does not depend on $\theta$), of size $c\varepsilon^3$; and $\abs{g_2}\le\abs{\hat f_0}+130E_2^2/(\bar\alpha_2r_*\sigma_3)\le c\varepsilon^3+c\,\varepsilon^5(\log\frac1\varepsilon)^{4\tau_0+2}/(\varepsilon\cdot\varepsilon\,K_2^{-\tau_0})$, which is the stated bound ($K_2\asymp\log\frac1\varepsilon$). The displacements: $\abs{\tilde\rho-\tilde\rho'}\le3E_2/(\bar\alpha_2\sigma_3)$, $\abs{J-J'}\le3E_2/(\bar\alpha_1\sigma_3)$, $\abs{\theta-\theta'}\le9E_2/(\bar\alpha_2r_*)$, with $\bar\alpha_2\asymp\varepsilon K_2^{-\tau_0}$.
\end{proof}

\subsection{Step E: the prepared system and its frequency map}\label{bpP:p11:sec-E}

Put $\tilde\rho_0=(2\pi\varepsilon,\dots,2\pi\varepsilon)$ and let $A_\varepsilon(\theta,r,\zeta,\rho')=(\theta,\ J_0+\varepsilon r,\ \zeta,\ \varepsilon(2\pi\mathbf1+\rho'))$ (here $\mathbf 1=(1,\dots,1)$). $A_\varepsilon$ is an $\varepsilon$-chart ($A_\varepsilon^*(dJ\wedge d\theta+d\tilde\rho\wedge d\zeta)=\varepsilon(dr\wedge d\theta+d\rho'\wedge d\zeta)$) commuting with all angle shifts. It maps the polydisc $\{\abs r_\infty<2,\ \abs{\rho'}_\infty<\frac\pi2\}$ into $\tfrac12V_2$ (radii $s_2/2$ for $J$ around $J_0$, $\pi\varepsilon/2$ for $\tilde\rho$ around $\tilde\rho_0$) if and only if $2\varepsilon\le s_2/2$ (the $\rho'$-radius $\frac\pi2$ is exactly the half polar radius). Define, with $\mathcal E=\{(\theta,r,\zeta,\rho'):\abs{\operatorname{Im}(\theta,\zeta)}<\sigma_3/6,\ \abs r_\infty<2,\ \abs{\rho'}_\infty<\frac\pi2\}$,
\[N_\varepsilon(r,\rho')=\tfrac1\varepsilon\bigl(h_0+g_2\bigr)(J_0+\varepsilon r,\ \varepsilon(2\pi\mathbf1+\rho')),\qquad P_\varepsilon=\tfrac1\varepsilon f_3\circ A_\varepsilon,\qquad\hat H_\varepsilon=\tfrac1\varepsilon H_2\circ A_\varepsilon=N_\varepsilon+P_\varepsilon ,\]
where $H_2:=H\circ\Phi^1_\varepsilon\circ\hat\Theta\circ\Psi_2$ is the Hamiltonian of Proposition~\ref{bpP:p11:avg-slow}(c). The compact set and the data of Definition~\ref{bpP:p11:kam-input} are
\[K_\varepsilon=[-\tfrac12,\tfrac12]^{2n},\ \ \vartheta=\tfrac14,\ \ \varepsilon^*=\tfrac12,\ \ s=s_*:=\sigma_3/6=\sigma_F/72,\ \ N=2n,\ \ \mu=\bar\mu,\ \ \tau_D=\max\{n\bar\mu,2n\},\]
so $\operatorname{meas}K_\varepsilon=1$, $d=1$, and the complex neighbourhoods $K_\varepsilon+2\vartheta\subset K_\varepsilon+3\vartheta\subset K_\varepsilon+4\vartheta\subset D(0,\frac32)^{2n}$ lie in $\{\abs r<2,\ \abs{\rho'}<\frac\pi2\}$. \emph{Domain table} (all in the $(J,\tilde\rho)$ variables, $\ell=\log\frac1\varepsilon$): $g_2$ and $f_3$ live on $\frac12V_2$, i.e.\ $\lvert J-J_0\rvert<s_2/2$, $\lvert\tilde\rho_j-2\pi\varepsilon\rvert<\pi\varepsilon/2$; the points used for estimates are those of $K_\varepsilon+3\vartheta$, i.e.\ $\lvert J-J_0\rvert<\frac54\varepsilon$ and $\lvert\tilde\rho_j-2\pi\varepsilon\rvert<\frac54\varepsilon$, so (for $\varepsilon<s_2/5$) the Cauchy discs of radius $s_2/4$ in $J$ and $\varepsilon/4$ in $\tilde\rho$ fit inside $\frac12V_2$ (the slack in $\tilde\rho$ is $(\frac\pi2-\frac54)\varepsilon>\frac14\varepsilon$); $\hat Z,\hat G_0$ live on $D(J_0,s_1/2)\times D(0,\pi r_z^2/16)$ with Cauchy discs of radius $s_1/4$ in $J$ and $\pi\varepsilon/4$, resp.\ $\pi r_z^2/32$, in $\tilde\rho$. Frequency map of $N_\varepsilon$: $\hat\omega_\varepsilon=\nabla_{(r,\rho')}N_\varepsilon=(\partial_Jh,\ \partial_{\tilde\rho}h)(J_0+\varepsilon r,\varepsilon(2\pi\mathbf1+\rho'))$ with $h=h_0+g_2$. It is convenient to write $\hat\omega_\varepsilon=(\hat\omega^{(1)},\varepsilon\hat\omega^{(2)})$ with
\[\hat\omega^{(1)}=\omega+\varepsilon\partial_J\hat G_0+\partial_Jg_2,\qquad\hat\omega^{(2)}=\tilde\Omega+\partial_{\tilde\rho}\hat Z+\varepsilon^{-1}\partial_{\tilde\rho}g_2\qquad\text{(arguments as above).}\]

\begin{proposition}[The prepared system]\label{bpP:p11:prepared}
Put $Q=5\tau_0+2+(\tau_0+1)(\bar\mu+1)$ and $\ell=\log\frac1\varepsilon$. There are $\varepsilon_3\in(0,\varepsilon_2]$ and $c_3\ge1$, depending only on the data and $\nu_2$, such that for $0<\varepsilon\le\varepsilon_3$: $5\varepsilon\le s_2$ (so $2\varepsilon\le s_2/2$ and the domain table of \S\ref{bpP:p11:sec-E} holds), $\theta_\varepsilon:=c_3\varepsilon\ell^Q\le1$, and the following hold. $N_\varepsilon,P_\varepsilon$ are holomorphic on $\mathcal E$, $1$-periodic in the angles, real on real points, $N_\varepsilon$ depends on $(r,\rho')$ only, and $H_2\circ A_\varepsilon=\varepsilon\hat H_\varepsilon$. Moreover:
\begin{enumerate}[label=(\alph*)]
\item $\sup_{\mathcal E}\lvert P_\varepsilon\rvert\le\varepsilon_f(\varepsilon):=c_3\,\varepsilon^{\nu_2+3/2}\ell^{2\tau_0+1}$;
\item $\lvert\hat\omega_\varepsilon\rvert_\infty\le C_1:=2L+1$ and $\lVert D\hat\omega_\varepsilon\rVert=\lvert D^2N_\varepsilon\rvert_{\infty\to\infty}\le C_1$ on $K_\varepsilon+3\vartheta$ (complex points);
\item (\emph{$J$-derivatives}) for $\lvert\alpha\rvert\le\bar\mu$ ($J$-multi-indices) and all real $J$ with $\lvert J-J_0\rvert_\infty<s_2/4$ and $\tilde\rho_j\in\varepsilon(2\pi+(-\frac54,\frac54))$: $\lvert\partial_J^\alpha(\hat\omega^{(i)}-\varpi^{(i)})\rvert\le\theta_\varepsilon$ ($i=1,2$; $\varpi^{(1)}=\omega$, $\varpi^{(2)}=\tilde\Omega$); in the derivative norm of Definition~\ref{bpP:p11a:index-def} the bound is $c(n,\bar\mu)\theta_\varepsilon$, which we absorb into $c_3$;
\item (\emph{full derivatives}) for $k\le\bar\mu+1$ on $K_\varepsilon+\vartheta$: $\lvert D^kg_2[\hat\omega_\varepsilon](c,y)\rvert_{\mathrm{op}}\le M_*:=2^{\bar\mu+1}(2L+1)^2$ ($c\in S^{2n-1}$), and $\lVert\hat\omega_\varepsilon\rVert^{\bar\mu}_{\mathrm{dg},K_\varepsilon+\vartheta}\le M_0^*:=c(n,\bar\mu)(2L+1)\ge1$, both independent of $\varepsilon$.
\end{enumerate}
\end{proposition}
\begin{proof}
\proofstep{Holomorphy, reality, periodicity}\label{stepP:p11:pr-basic} $\Phi^1_\varepsilon$, $\hat\Theta$, $\Psi_2$ are real holomorphic with the lattice properties and $H_2\circ A_\varepsilon=H\circ(\Phi^1_\varepsilon\circ\hat\Theta\circ\Psi_2\circ A_\varepsilon)$; the domain chain is the following: $A_\varepsilon(\mathcal E)\subset\T_{\sigma_3/6}\times\frac12V_2$ (domain table, $2\varepsilon\le s_2/2$), $\Psi_2$ maps it into $\T_{\sigma_3}\times V_2$, on which $\hat\Theta$ is defined with values in $\T_{\sigma_F/6}\times D(J_0,s_2)\times B_{r_z/4}\subset\mathcal D_1'$ (the action $J$ is unchanged, so lies in $D(J_0,s_2)\subset D(J_0,s_1/2)$; the angular strip may be taken $\T_{\sigma_3}$) (here $\sigma_3=\sigma_F/12$, $s_2\le s_1/4$, Prop.~\ref{bpP:p11:avg-slow} and $\Pi(\mathcal P_\varepsilon)\subset B_{r_z/4}$). $N_\varepsilon=\frac1\varepsilon(h_0+g_2)\circ A_\varepsilon$ is a function of $(r,\rho')$, $P_\varepsilon=\varepsilon^{-1}f_3\circ A_\varepsilon$.
\proofstep{(a)}\label{stepP:p11:pr-a} $\sup\lvert f_3\rvert\le\lVert f_3\rVert\le\varepsilon^{\nu_2}E_2$ (\Nref{bpN:p4:sup-weighted}(a)); divide by $\varepsilon$.
\proofstep{Derivative inputs}\label{stepP:p11:pr-in} Write $h=h_0+g_2$ and use $s_1\ge c\ell^{-(\tau_0+1)}$, $s_2\ge c\ell^{-(\tau_0+1)}$ (from \eqref{bpP:p11:eq:par1}, \eqref{bpP:p11:eq:par2}, $K_1,K_2\asymp\ell$). (i) $\lvert\partial_J^\beta\hat G_0\rvert\le c\,M_F\beta!(4/s_1)^{\lvert\beta\rvert}\le c\,\ell^{(\tau_0+1)\lvert\beta\rvert}$ (e2, Cauchy radius $s_1/4$); (ii) $\lvert\partial_J^\alpha\partial_{\tilde\rho}\hat Z\rvert\le c_2(\varepsilon+\eta_1)\alpha!(4/s_1)^{\lvert\alpha\rvert}\le c\,\varepsilon\ell^{2\tau_0+1+(\tau_0+1)\lvert\alpha\rvert}$ (Step~\ref{stepP:p11:as-osc}); (iii) $\lvert g_2\rvert\le c_2\varepsilon^3\ell^{5\tau_0+2}$ on $\frac12V_2$, hence by (e2) with radii $s_2/4$ in $J$ and $\varepsilon/4$ in $\tilde\rho$ (they fit: domain table): $\lvert\partial_J^\alpha\partial_{\tilde\rho}^\gamma g_2\rvert\le c\,\alpha!\gamma!\,\varepsilon^{3-\lvert\gamma\rvert}\ell^{5\tau_0+2+(\tau_0+1)\lvert\alpha\rvert}$.
\proofstep{(b)}\label{stepP:p11:pr-b} $\partial_rN_\varepsilon=\partial_Jh=\omega+\varepsilon\partial_J\hat G_0+\partial_Jg_2$ and $\partial_{\rho'}N_\varepsilon=\partial_{\tilde\rho}h=\varepsilon(\tilde\Omega+\partial_{\tilde\rho}\hat Z)+\partial_{\tilde\rho}g_2$. By Step~\ref{stepP:p11:pr-in} (and Proposition~\ref{bpP:p11:holo-chart}(a) for $\omega,\tilde\Omega$) these are $\le L+c\varepsilon\ell^{\tau_0+1}+c\varepsilon^3\ell^{Q}$ and $\le\varepsilon L+c\varepsilon^2\ell^{2\tau_0+1}+c\varepsilon^2\ell^{5\tau_0+2}$; all corrections are $\le1$ for $\varepsilon\le\varepsilon_3$ (e4), so $\lvert\hat\omega_\varepsilon\rvert_\infty\le2L+1$. Hessian: $\partial_r^2N_\varepsilon=\varepsilon\partial_J^2h$, $\partial_r\partial_{\rho'}N_\varepsilon=\varepsilon\partial_J\partial_{\tilde\rho}h$, $\partial_{\rho'}^2N_\varepsilon=\varepsilon\partial_{\tilde\rho}^2h$, with $\partial_{\tilde\rho}^2h=\varepsilon\partial_{\tilde\rho}^2\hat Z+\partial^2_{\tilde\rho}g_2$; every entry is $\le c\varepsilon^{a}\ell^{b}$ with $a\ge1$ by Step~\ref{stepP:p11:pr-in} (for $\partial^2_{\tilde\rho}\hat Z$ use (e2) for $\hat Z$ with radius $\pi r_z^2/32$), so the operator norm is $\le2n\cdot c\varepsilon^a\ell^b\le1\le C_1$ for $\varepsilon\le\varepsilon_3$.
\proofstep{(c)}\label{stepP:p11:pr-c} $\hat\omega^{(1)}-\omega=\varepsilon\partial_J\hat G_0+\partial_Jg_2$ and $\hat\omega^{(2)}-\tilde\Omega=\partial_{\tilde\rho}\hat Z+\varepsilon^{-1}\partial_{\tilde\rho}g_2$. For $\lvert\alpha\rvert\le\bar\mu$, Step~\ref{stepP:p11:pr-in} gives: $\varepsilon\lvert\partial_J^{\alpha+1}\hat G_0\rvert\le c\varepsilon\ell^{(\tau_0+1)(\bar\mu+1)}$; $\lvert\partial_J^{\alpha+1}g_2\rvert\le c\varepsilon^3\ell^{5\tau_0+2+(\tau_0+1)(\bar\mu+1)}=c\varepsilon^3\ell^Q$; $\lvert\partial_J^\alpha\partial_{\tilde\rho}\hat Z\rvert\le c\varepsilon\ell^{2\tau_0+1+(\tau_0+1)\bar\mu}$; $\varepsilon^{-1}\lvert\partial_J^\alpha\partial_{\tilde\rho}g_2\rvert\le c\varepsilon\ell^{5\tau_0+2+(\tau_0+1)\bar\mu}$. All exponents are $\le Q$ and all powers of $\varepsilon$ are $\ge1$, so each is $\le c\varepsilon\ell^Q$. The extra factor $(4/s)^{\lvert\alpha\rvert}\asymp\ell^{(\tau_0+1)\lvert\alpha\rvert}$ per $J$-derivative is the loss that a bound by the common power $5\tau_0+2$ would omit (that bound is false for, e.g., $g_2=\varepsilon^3\ell^{5\tau_0+2}\frac{J-J_0}{s_2}\frac{\tilde\rho-2\pi\varepsilon}\varepsilon$ of the stated sup size); summing the four terms and the finite number of orders gives $c_3\varepsilon\ell^Q$.
\proofstep{(d)}\label{stepP:p11:pr-d} By Lemma~\ref{bpP:p11a:M-bound}, $\lvert D^kg_2[\hat\omega_\varepsilon]\rvert_{\mathrm{op}}\le2^kM_1^2$ with $M_1=\lVert\hat\omega_\varepsilon\rVert^{\bar\mu+1}_{\mathrm{op},K_\varepsilon+\vartheta}$. The main terms of $\hat\omega_\varepsilon$, $(\omega(J_0+\varepsilon r),\varepsilon\tilde\Omega(J_0+\varepsilon r))$, are independent of $\rho'$, and each $r$-derivative of order $j\ge1$ carries a factor $\varepsilon^j$, so their norms are $\le L$ (Proposition~\ref{bpP:p11:holo-chart}(a)). In the remaining terms an $(r,\rho')$-derivative of order $j=j_J+j_\rho$ is $\varepsilon^{j}$ times a derivative in $(J,\tilde\rho)$ of order $(j_J,j_\rho)$, and by Step~\ref{stepP:p11:pr-in} each such term is $\le c\varepsilon^{j}\cdot\varepsilon^{1-j_\rho}\ell^{b}\le c\varepsilon\ell^b$ (the factor $\varepsilon^{-j_\rho}$ of the $\tilde\rho$-Cauchy radius is cancelled by $\varepsilon^{j_\rho}$ from the rescaling; the $\hat Z$ terms use the fixed radius and are smaller), with $b$ a fixed number depending on $\bar\mu,\tau_0$; hence $\le1$ for $\varepsilon\le\varepsilon_3$. So $M_1\le2L+1$ and $M_*$ is as claimed. The diagonal norm $\lVert\cdot\rVert_{\mathrm{dg}}$ (Euclidean unit sphere) is bounded by $c(n,\bar\mu)$ times the max-norm operator norm, giving $M_0^*$.
\end{proof}

\begin{lemma}[Directional stability and target scaling; R-CP-2]\label{bpP:p11:amount}
Let $L_1:=c(n,\bar\mu)L$ be a bound for $\lVert\varpi_J\rVert^{\bar\mu}_{\mathrm{dg}}$ on $\mathcal J\cap\R^n$ (Proposition~\ref{bpP:p11:holo-chart}(a)), and decrease $\varepsilon_3$ (depending only on the data and the already fixed $\nu_2$) so that for $0<\varepsilon\le\varepsilon_3$ also $2^{\bar\mu}\theta_\varepsilon(2L_1+\theta_\varepsilon)\le\frac12\beta_J$ and $\beta_\varepsilon\le\min\{1,M_*\}$, with $\theta_\varepsilon$ of Proposition~\ref{bpP:p11:prepared}(c) and $\beta_J$ of Lemma~\ref{bpP:p11:norm-dict}. Then, in the sense of Definition~\ref{bpP:p11a:index-def} with $\rho=2$, $m=N=2n$,
\[\beta_{2,\bar\mu}(\hat\omega_\varepsilon,K_\varepsilon)\ \ge\ \beta_\varepsilon:=\tfrac12\,\varepsilon^{\bar\mu+2}\beta_J>0 .\]
\end{lemma}
\begin{proof}
\proofstep{Normalisation of the second block}\label{stepP:p11:am-norm} $\hat\omega_\varepsilon=\operatorname{diag}(\mathbf1_n,\varepsilon\mathbf1_n)\check\omega_\varepsilon$, $\check\omega_\varepsilon=(\hat\omega^{(1)},\hat\omega^{(2)})$. For $D=\operatorname{diag}(\mathbf1,\varepsilon\mathbf1)$ and $c\in S^{2n-1}$: $\langle c,D\check\omega\rangle=\langle Dc,\check\omega\rangle=\lvert Dc\rvert\langle\hat c,\check\omega\rangle$, $\hat c=Dc/\lvert Dc\rvert\in S^{2n-1}$, $\varepsilon\le\lvert Dc\rvert\le1$; so $g_2[D\check\omega](c,\cdot)=\lvert Dc\rvert^2g_2[\check\omega](\hat c,\cdot)$ and $\beta_{2,\mu}(D\check\omega,K)\ge\varepsilon^2\beta_{2,\mu}(\check\omega,K)$.
\proofstep{Freeze $\rho'$; stability in $J$}\label{stepP:p11:am-dir} Fix $\rho'_1\in[-\frac12,\frac12]^n$, $\tilde\rho_1=\varepsilon(2\pi\mathbf1+\rho'_1)$, and on $V=D_\R(J_0,s_2/4)\subset\R^n$ put $\psi_\varepsilon(J)=\check\omega_\varepsilon$ evaluated at $(J,\tilde\rho_1)$ and $\psi=\varpi_J$; both are $C^{\bar\mu}$ maps $V\to\R^{2n}$. Let $A=\{J_0+\varepsilon r:\lvert r\rvert_\infty\le\frac12\}$, a compact subset of $V$ since $\varepsilon/2<s_2/4$ ($5\varepsilon\le s_2$), and $A\subset K_J$ for small $\varepsilon$ ($J_0\in\operatorname{int}K_J$). Proposition~\ref{bpP:p11:prepared}(c) (its points $\tilde\rho_1$ and $J\in V$ are in its range) gives $\theta:=\lVert\psi_\varepsilon-\psi\rVert^{\bar\mu}_{\mathrm{dg},A}\le\theta_\varepsilon$, and $\lVert\psi\rVert^{\bar\mu}_{\mathrm{dg},A}\le L_1$. Lemma~\ref{bpP:p11a:stability} (with $N=n$, $m=2n$, $\rho=2$, $M_0=L_1$) and $\beta_{2,\bar\mu}(\psi,A)\ge\beta_J$ (Lemma~\ref{bpP:p11:norm-dict}), $2^{\bar\mu}\theta(2L_1+\theta)\le\frac12\beta_J$, give $\beta_{2,\bar\mu}(\psi_\varepsilon,A)\ge\frac12\beta_J$, the $n$-dimensional amount being the minimum over $(c,J)$ of $\max_{\nu\le\bar\mu}\max_{a\in S^{n-1}}\lvert D_J^\nu g_2[\psi_\varepsilon](c,J)(a^\nu)\rvert$.
\proofstep{Rescaling of the actions}\label{stepP:p11:am-scale} In the variables $(r,\rho')$ the derivative of $g_2[\check\omega_\varepsilon](c,\cdot)$ of order $\nu$ in a direction $(a,0)$, $a\in S^{n-1}$, at $(r,\rho'_1)$ equals $\varepsilon^\nu D_J^\nu g_2[\psi_\varepsilon](c,J)(a^\nu)$, $J=J_0+\varepsilon r$ (chain rule). The full tensor norm (maximum over all unit directions of $\R^{2n}$, Definition~\ref{bpP:p11a:index-def}) is at least its value on directions $(a,0)$. Since $\max_{\nu\le\bar\mu}\varepsilon^\nu x_\nu\ge\varepsilon^{\bar\mu}\max_\nu x_\nu$ ($\varepsilon<1$): $\beta_{2,\bar\mu}(\check\omega_\varepsilon,K_\varepsilon)\ge\varepsilon^{\bar\mu}\frac12\beta_J$, uniformly in $\rho'_1$. With Step~\ref{stepP:p11:am-norm}: $\beta_{2,\bar\mu}(\hat\omega_\varepsilon,K_\varepsilon)\ge\varepsilon^{\bar\mu+2}\frac12\beta_J$.
\end{proof}

\begin{remark}[R-CP-2 and what is not claimed]\label{bpP:p11:r-cp-2}
CP09 p.~17 asserts $\beta_\varepsilon\ge\varepsilon^{\bar\mu}\bar\beta/2$ from $\Psi_\varepsilon=\Psi_0+O(\varepsilon)$ in $C^0$ of the \emph{rescaled} variable; that is insufficient (the amount is $\asymp\varepsilon^{\bar\mu}$, a $C^0$-perturbation of size $\varepsilon$ is larger). Here the perturbation is measured in $C^{\bar\mu}$ in the \emph{unscaled} action $J$ (Proposition~\ref{bpP:p11:prepared}(c)), where it is $O(\varepsilon\log^{Q}\frac1\varepsilon)$ with $Q=5\tau_0+2+(\tau_0+1)(\bar\mu+1)$ (Proposition~\ref{bpP:p11:prepared}(c); after decreasing the threshold, in particular $O(\varepsilon^{1/2})$) against the unscaled amount $\beta_J$; the rescaling then multiplies both by the same $\varepsilon^\nu$ (Lemma~\ref{bpP:p11:amount}). The $\rho'$-dependence of $N_\varepsilon$ is not controlled in $C^{\bar\mu}$ at all (only through Proposition~\ref{bpP:p11:prepared}(d), a bound), and is not needed: the amount in the full sense is bounded below by the amount in the $r$-directions. Also corrected: CP09 (17), (28) and Prop.~2 carry a spurious factor $\frac12$ in front of $\hat\Omega\cdot(\rho_0+\varepsilon\rho)$ (with $z_j=\sqrt{2(\rho_{0j}+\rho_j)}e^{-i\zeta_j}$ one has $\frac12(u_j^2+v_j^2)=\rho_{0j}+\rho_j$), which changes only the constant in the amount.
\end{remark}

\subsection{Step F: the local wrapper and the transport of the tori}\label{bpP:p11:sec-F}

\begin{lemma}[Forward tube for $n$ planets]\label{bpP:p11:forward-tube}
Let $a_i>0$, $\Lambda_i^*=m_i\sqrt{\kappa a_i}$, $\zeta>0$. There is $\delta_\zeta\in(0,\frac14\min_i\Lambda_i^*]$, depending only on $(G,m_0,m,a,\zeta)$, such that every $w=((\lambda,\xi),(\Lambda,\eta))\in\R^{4n}$ with $\abs{\Lambda_i-\Lambda_i^*}\le\delta_\zeta$ and $\Gamma_i(w)\le\delta_\zeta$ for all $i$ lies in $\mathcal O_n$ and satisfies $\abs{\Psi^{(n)}(w)-\operatorname{circ}(\lambda/2\pi)}_\infty<\zeta$; in particular $\Psi^{(n)}(w)\in U_\zeta$.
\end{lemma}
\begin{proof}
Steps~\ref{stepP:p12:ft-cpt}--\ref{stepP:p12:ft-cmp} of Lemma~\ref{bpP:p12:forward-tube} with $2$ replaced by $n$ in the index ranges and $\R^8$ by $\R^{4n}$; nothing else changes (compact set $Q_0$, uniform continuity (e3) of $\Psi^{(n)}$, $\Psi^{(n)}(w_0)=\operatorname{circ}(\lambda/2\pi)$ by Lemma~\ref{bpP:p7:tube}).
\end{proof}

\begin{theorem}[Local wrapper for P-FIXED: tori in the tube]\label{bpP:p11:wrapper}
Let $n\ge2$, $G,m_0>0$, $m_i>0$, $0<a_1<\dots<a_n$, $\eta>0$, and $\zeta_*,d,D$ as in Definition~\ref{bpP:p6:tube-constants}. Then there is $\varepsilon_0>0$ such that for every $\varepsilon\in(0,\varepsilon_0)$ there is a Borel set $S_\varepsilon\subset\Ph{2n}$ with $0<\operatorname{meas}S_\varepsilon<\infty$, $S_\varepsilon\subset U_{\zeta_*}$, and
\begin{multline*}\forall z\in S_\varepsilon\ \ \exists\,w\in\R^{2n},\ \mathcal K:\R^{2n}\to\Ph{2n},\ \theta_0\in\R^{2n}:\\ z=\mathcal K(\theta_0)\ \wedge\ \mathcal K(\R^{2n})\subset S_\varepsilon\ \wedge\ \operatorname{IsPhysTorus}(H_{\rm fix},w,\mathcal K),\end{multline*}
for $H_{\rm fix}=\texttt{fixedHamiltonian}\ G\ m_0\ m\ \varepsilon$. (The frequency $w$ depends on the point, as in Definition~\ref{bpP:p0:fixkol}; Part~13 applies Proposition~\ref{bpP:p6:fixkol-assembly} to obtain $\operatorname{FixKol}(G,m_0,m,a,\varepsilon,\eta,d,D,S_\varepsilon)$ and the frozen P-FIXED.) The order of choice is: data $\to$ patch, $(\bar\mu,\bar\beta)$, $J_0,\gamma_0$ $\to$ $r_C,\sigma_F,r_z,M_F,L,B_\varrho,B'_\varrho$ $\to$ the fixed part $\mathfrak d_*$ of the KAM datum and its constants $B_0=\bar\mu(\bar\mu+1)+2$, $c_*,c_D,\kappa_D,a_D$ $\to$ $\nu_2$ $\to$ $\varepsilon_0$ $\to$ $\varepsilon$. The conclusions (K2)--(K4) and the displacement power (K6) enter as the cited input of Definition~\ref{bpP:p11:kam-input}; they are \emph{not} proved in this part.
\end{theorem}

\begin{remark}[The data of the KAM input]\label{bpP:p11:data-star}
For $0<\varepsilon\le\varepsilon_3$ the Hamiltonian $\hat H_\varepsilon=N_\varepsilon+P_\varepsilon$ on $\T_{s_*}\times(K_\varepsilon+2\vartheta)$ is $\mathfrak d_\varepsilon$-admissible with $\mathfrak d_\varepsilon=(\bar\mu,\beta_\varepsilon,K_\varepsilon,\frac14,s_*,M_*,M_0^*,C_1,\tau_D,\frac12)$: holomorphy, periodicity, reality: Proposition~\ref{bpP:p11:prepared}; (A1): Lemma~\ref{bpP:p11:amount}; (A2): Proposition~\ref{bpP:p11:prepared}(d); (A3): (b), $C_1=2L+1\ge3$; (A4): (d) with $M_0^*$; $0<\beta_\varepsilon\le\min\{1,M_*\}$ by the choice of $\varepsilon_3$ (note $M_*\ge\beta_\varepsilon$ as required). The entries are: $K_\varepsilon=[-\frac12,\frac12]^{2n}$, $d_\varepsilon=\operatorname{diam}_\infty K_\varepsilon=1$ (this is the diameter of the KAM datum, not the physical lower distance $d$ of Definition~\ref{bpP:p6:tube-constants}), $\vartheta=\frac14$, $s_*=\sigma_F/72$, $\varepsilon^*=\frac12$, $N=2n$, $\mu=\bar\mu$, $\tau_D=\max\{n\bar\mu,2n\}>\max\{n\bar\mu-1,2n-1\}$, and $M_*,M_0^*,C_1$ independent of $\varepsilon$ (they depend on $L$ only). Hence all entries except $\beta_\varepsilon=\frac12\beta_J\varepsilon^{\bar\mu+2}$ form the $\varepsilon$-independent part $\mathfrak d_*$; (K5) and (K6) apply to this family with the constants of $\mathfrak d_*$, and $\varepsilon_{\rm adm}(\mathfrak d_\varepsilon)\ge c_*\beta_\varepsilon^{B_0}$, $B_0=\bar\mu(\bar\mu+1)+2$.
\end{remark}

\begin{proof}[Proof of Theorem~\ref{bpP:p11:wrapper}]
\proofstep{Choice of $\nu_2$ and $\varepsilon_0$}\label{stepP:p11:wr-nu} Let $c_*,c_D,\kappa_D,a_D$ be the constants of (K5), (K6) for $\mathfrak d_*$ and $B_0=\bar\mu(\bar\mu+1)+2$. Choose $\nu_2:=\max\bigl\{(\bar\mu+2)B_0+2,\ \tfrac{(\bar\mu+2)a_D+1}{\kappa_D}+2\bigr\}$; it depends on no $\varepsilon$ and no constant of Steps~C--E (the fixed datum and its (K5), (K6) constants are independent of $\nu_2$; the intermediate estimate constants and thresholds of Steps~C--E may depend on the subsequently fixed $\nu_2$). By Proposition~\ref{bpP:p11:prepared}(a), $\varepsilon_f(\varepsilon)\le c_3\varepsilon^{\nu_2+3/2}\ell^{2\tau_0+1}$, while $\varepsilon_{\rm adm}(\mathfrak d_\varepsilon)\ge c_*(\tfrac12\beta_J)^{B_0}\varepsilon^{(\bar\mu+2)B_0}\ge c'\varepsilon^{\nu_2-2}$ and $\Delta\le c_D\varepsilon_f^{\kappa_D}(\tfrac12\beta_J\varepsilon^{\bar\mu+2})^{-a_D}\le c''\varepsilon^{\kappa_D(\nu_2+3/2)-(\bar\mu+2)a_D}\ell^{c}\le c''\varepsilon^{1+2\kappa_D}\ell^c$. Hence there is $\varepsilon_0\le\varepsilon_3$ (depending only on the data, $\nu_2$ and the constants above) such that for every $\varepsilon\in(0,\varepsilon_0)$: $\varepsilon_f(\varepsilon)\le\varepsilon_{\rm adm}(\mathfrak d_\varepsilon)$, $\Delta\le\frac12$, and the tube inequalities of Step~\ref{stepP:p11:wr-tube} hold. This is the quantifier order demanded by R-CP-1 and A4: the admissible size is bounded \emph{below} by a power, and $\nu_2$ is fixed \emph{before} $\varepsilon$.
\proofstep{The KAM input}\label{stepP:p11:wr-kam} Apply Definition~\ref{bpP:p11:kam-input} to $\hat H_\varepsilon$ (Remark~\ref{bpP:p11:data-star}) with $\varepsilon_f=\varepsilon_f(\varepsilon)$: a compact $\mathcal H_\varepsilon\subset K_\varepsilon$ with $\operatorname{meas}\mathcal H_\varepsilon\ge\frac12$, and for $b\in\mathcal H_\varepsilon$ maps $\mathcal K_b=(\theta+u^\theta_b,\ \zeta+u^\zeta_b;\ b_r+u^r_b,\ b_{\rho'}+u^{\rho'}_b)$ (positions $(\theta,\zeta)$, actions $(r,\rho')$) with defect $\lvert u_b\rvert_\infty\le\Delta\le\frac12$; hence the actions of every real point of every torus satisfy $\lvert r\rvert_\infty\le1$, $\lvert\rho'\rvert_\infty\le1<\frac\pi2$.
\proofstep{Chain of charts}\label{stepP:p11:wr-chain} $\Xi_\varepsilon:=\Psi_c\circ\Phi^1_\varepsilon\circ\hat\Theta\circ\Psi_2\circ A_\varepsilon$ is defined and real analytic on $\mathcal R_\varepsilon=\{(\theta,r,\zeta,\rho')\in\R^{4n}:\lvert r\rvert_\infty<2,\ \lvert\rho'\rvert_\infty<\frac\pi2\}$ (domain chain of Step~\ref{stepP:p11:pr-basic}; $\Psi_c$ on real points of $\mathcal O_c$: $\lvert\Pi\rvert\le r_z/4$, $J$ real in $D(J_0,s_1)$). All charts are symplectic for the single form $dJ\wedge d\theta+dy\wedge dx$ on the working side, which on the polar side is $dJ\wedge d\theta+d\tilde\rho\wedge d\zeta$ (Lemma~\ref{bpP:p11:total-chart}) and on the rescaled side $\varepsilon(dr\wedge d\theta+d\rho'\wedge d\zeta)$ ($A_\varepsilon$): $\Xi_\varepsilon$ is an $\varepsilon$-chart in the sense of Definition~\ref{bpP:p1:chart}. Its lattice properties (Q1)--(Q2) with $\sigma_k=0$ and $\iota_k=((k,0),0)$ shifting $(\theta,\zeta)$ by $k\in\Z^{2n}$ follow factor by factor: $A_\varepsilon$ leaves the angles unchanged; $\Psi_2$ is equivariant and injective (Lemma~\ref{bpP:p11:aniso-avg}(a)); $\Phi^1_\varepsilon-\mathrm{id}$ is $\Z^n$-periodic in $\theta$ and $\Phi^1_\varepsilon$ is injective; $\hat\Theta$ and $\Psi_c$ by Lemma~\ref{bpP:p11:total-chart}. Moreover $H_{\rm fix}\circ\Xi_\varepsilon=H_2\circ A_\varepsilon=\varepsilon\hat H_\varepsilon$ on $\mathcal R_\varepsilon$ (Proposition~\ref{bpP:p11:holo-chart}(b), Propositions~\ref{bpP:p11:avg-fast}, \ref{bpP:p11:avg-slow}(c)), i.e.\ $\hat H_\varepsilon=H_{\rm fix}^{\Xi_\varepsilon}$ (Definition~\ref{bpP:p1:chart}); the physical time is preserved (Lemma~\ref{bpP:p1:chart-field}).
\proofstep{Physical tori, per point}\label{stepP:p11:wr-tor} Let $\mathcal B=\mathcal K(\mathcal H_\varepsilon\times[0,1)^{2n})$, $A=\mathcal K(\mathcal H_\varepsilon\times\R^{2n})$ and $S_\varepsilon:=\Xi_\varepsilon(A)$. For $b\in\mathcal H_\varepsilon$ the map $\mathcal K_b$ takes values in $\mathcal R_\varepsilon$ (Step~\ref{stepP:p11:wr-kam}) and satisfies (L1)--(L6) of Definition~\ref{bpP:p1:lattice-torus} with $G=\hat H_\varepsilon$ (differentiable: holomorphic), $O=\mathcal R_\varepsilon$, $w=w_b$ and $\iota_k=((k,0),0)$ (by (K2): $\mathcal K_b(\theta+k)=\mathcal K_b(\theta)+(k,0)$); $w_b$ is $(\gamma,\tau_D)$-Diophantine with $\tau_D\ge2n>2n-1$. Theorem~\ref{bpP:p1:torus-transport}(b) gives $\operatorname{IsPhysTorus}(H_{\rm fix},w_b,\Xi_\varepsilon\circ\mathcal K_b)$. Now let $z\in S_\varepsilon$: by definition $z=\Xi_\varepsilon(\mathcal K_b(\theta_0))$ for some $b\in\mathcal H_\varepsilon$, $\theta_0$; put $w:=w_b$ (a frequency chosen for this point), $\mathcal K:=\Xi_\varepsilon\circ\mathcal K_b$. Then $z=\mathcal K(\theta_0)$ and $\mathcal K(\R^{2n})=\Xi_\varepsilon(\mathcal K_b(\R^{2n}))\subset\Xi_\varepsilon(A)=S_\varepsilon$: the per-point statement.
\proofstep{Borel, positive measure}\label{stepP:p11:wr-meas} By (K4) $\mathcal B$ is Borel with $\operatorname{meas}\mathcal B>0$, and $A=\bigcup_{k\in\Z^{2n}}(\mathcal B+\iota_k)$ (equivariance of $\mathcal K_b$), a countable union of Borel sets, so $A$ is Borel, invariant under all $\iota_k$, and $A\subset\mathcal R_\varepsilon$ (open, invariant). Let $Q=\{(\theta,r,\zeta,\rho'):(\theta,\zeta)\in[0,1)^{2n}\}$ be the ambient cell. The translates $Q+\iota_k$ partition the ambient space, so $\mathcal B=\bigcup_k\mathcal B\cap(Q+\iota_k)$ and some $k$ has $\operatorname{meas}(\mathcal B\cap(Q+\iota_k))>0$; translating back, $\operatorname{meas}(A\cap Q)\ge\operatorname{meas}((\mathcal B-\iota_k)\cap Q)>0$ (the parameter cell $\mathcal B$ is \emph{not} identified with $A\cap Q$). By Lemma~\ref{bpP:p6:measure-positive}(b), applied to the $\varepsilon$-chart $\Xi_\varepsilon$ with (Q1)--(Q2), $S_\varepsilon=\Xi_\varepsilon(A)$ is Borel and $\operatorname{meas}S_\varepsilon=\varepsilon^{2n}\operatorname{meas}(A\cap Q)>0$ (the conformal Jacobian is $\varepsilon^{2n}$, Lemma~\ref{bpP:p1:liouville-lattice}); finiteness follows from the tube inclusion below ($S_\varepsilon\subset U_{\zeta_*}$, Lemma~\ref{bpP:p6:measure-positive}(a)).
\proofstep{Inclusion in the tube}\label{stepP:p11:wr-tube} Let $(\theta,r,\zeta,\rho')\in A$: $\lvert r\rvert_\infty\le1$, $\lvert\rho'\rvert_\infty\le1$ (Step~\ref{stepP:p11:wr-kam}); the real action coordinates are $J=J_0+\varepsilon r$ and $\tilde\rho_j=\varepsilon(2\pi+\rho'_j)\in\varepsilon[2\pi-1,2\pi+1]$. By Proposition~\ref{bpP:p11:avg-slow}(d) the map $\Psi_2$ changes $\tilde\rho$ by at most $c_2\varepsilon^{3/2}\ell^{3\tau_0+1}=o(\varepsilon)$ and $J$ by at most $c_2\varepsilon^{5/2}\ell^{3\tau_0+1}=o(\varepsilon)$ (these are the actual bounds of (d); both are $o(\varepsilon)$ for $\varepsilon\to0$ by (e4)), so after $\Psi_2$: $\tilde\rho_j\in\varepsilon[2\pi-2,2\pi+2]$ and $\lvert J-J_0\rvert_\infty\le2\varepsilon$ for $\varepsilon\le\varepsilon_0$. $\Pi$ gives $x_j^2+y_j^2=\tilde\rho_j/\pi\le(2\pi+2)\varepsilon/\pi<3\varepsilon$; $\Phi^1_\varepsilon$ moves $J$ by $\le3\varepsilon M_F/(2\pi\alpha_1\sigma_F)\le c\varepsilon\ell^{\tau_0}$ and $z$ by $\le9\varepsilon M_F/(2\pi\alpha_1r_z)\le c\varepsilon\ell^{\tau_0}=o(\sqrt\varepsilon)$ (Proposition~\ref{bpP:p11:avg-fast}). Finally $(\xi,\eta)=\varrho(x,y)$ with $\varrho(\Lambda)\in SO(n)$ \emph{real orthogonal at these real points} (Proposition~\ref{bpP:p9:lift}), so $\sum_i(\xi_i^2+\eta_i^2)=\sum_j(x_j^2+y_j^2)\le4n\varepsilon$ and $\Gamma_i=\frac12(\xi_i^2+\eta_i^2)\le2n\varepsilon$; and $\lvert\Lambda-\Lambda^*\rvert_\infty\le\lvert\Lambda_0-\Lambda^*\rvert+c\varepsilon\ell^{\tau_0}<\delta_*/8+c\varepsilon\ell^{\tau_0}\le\delta_*$ (patch choice). For $\varepsilon\le\varepsilon_0$ the hypotheses of Lemma~\ref{bpP:p11:forward-tube} (with $\zeta=\zeta_*$, $\delta_{\zeta_*}=\delta_*$) hold, so $\Xi_\varepsilon(\theta,r,\zeta,\rho')=\Psi^{(n)}(\cdots)\in U_{\zeta_*}$: $S_\varepsilon\subset U_{\zeta_*}\subset U_\eta$ (Lemma~\ref{bpP:p6:tube-geometry}(a)).
\end{proof}

\begin{remark}[Displacement in the polar and rotation charts]\label{bpP:p11:displacement}
In the polar chart the oscillator actions $\tilde\rho_j$ of every torus lie in $\varepsilon[2\pi-2,2\pi+2]$ (Step~\ref{stepP:p11:wr-tube}), in particular they are positive, uniformly over the family; this is the quantitative form of CP09 Theorem~5 ($\bar w$ of norm $2\varepsilon$). There is no Birkhoff chart in this route (it is replaced by Step D); the rotation chart $\Phi_\varrho$ preserves $\Lambda$ and, at real points, the Euclidean norm of $(\xi,\eta)$, which is what the tube needs.
\end{remark}

\subsection{Sources, deviations, errata, forward references}\label{bpP:p11:sec-notes}

\paragraph{Sources read.} CP09 (key CP09 in \texttt{docs/sources.json}; arXiv v1 numbering, 26 pages) pp.~5--8 (Def.~1, Lemma~1, Def.~2, Thm.~3, Thm.~4 and Remark~4), pp.~9--17 (\S2.3.1: Prop.~1, Lemmas~2--7, (11)--(30); \S2.3.2: Prop.~2), pp.~17--21 (\S2.3.3: Lemma~8, Remark~5, (34)--(47)); \texttt{reviews/p11-route-study.md}, \texttt{reviews/p11a-russmann-check.md} (R-CP-1, R-CP-2, A2--A7), \texttt{reviews/p11-review-1.md} (F01--F14); Parts~1, 2, 6--10, 11a-1, 11a-2, 11a-3 (\texttt{adm-budget}, \texttt{adm-budget-beta}) of this blueprint; \texttt{PLAN-11a.md} \S4; OVERNIGHT.md entries of 21:15 (Tracks~I, N, P).

\paragraph{Deviations.} \textbf{D-11-1} Steps~C--D replace CP09 Lemmas 4--7: fast averaging as in CP09 Lemma 4 (derived from Part~2), then \emph{no} equilibrium shift, quadratic diagonalisation or finite Birkhoff form; instead a polar normalisation on the shell $\tilde\rho\approx2\pi\varepsilon$ and a second two-scale averaging over all angles with two action radii (Lemma~\ref{bpP:p11:aniso-avg}; needs $\tilde\Omega(J_0)$ Diophantine, Lemma~\ref{bpP:p11:select}). \textbf{D-11-2} $\nu_1=3$ is fixed; the arbitrary order is carried by $\nu_2$ (exponential smallness in the slow averaging), not by $\nu_1$. \textbf{D-11-3} The rescaled data are $\vartheta=\frac14$, $\varepsilon^*=\frac12$, $K_\varepsilon=[-\frac12,\frac12]^{2n}$, independent of $\varepsilon$; CP09 uses the smaller sets $\vartheta=\varepsilon/16$, $\varepsilon^*=\varepsilon^{p+1}$ inside the larger pulled-back domain, which is legitimate; we use the whole half polar domain instead, which makes every datum but the amount independent of $\varepsilon$. \textbf{D-11-4} Max norms and period one throughout; the Euclidean balls of CP09 (11), (40) are replaced by polydiscs (Steps~\ref{stepP:p11:hc-lift}, \ref{stepP:p11:wr-tube}). \textbf{D-11-5} Part~11 uses the amount only through the $r$-directions (Lemma~\ref{bpP:p11:amount}); the $\rho'$-dependence of $N_\varepsilon$ is only bounded.

\paragraph{Suspected errata (for the ledger).} E-P-CP09-1: (17), (28), Prop.~2: factor $\frac12$ in front of $\hat\Omega\cdot(\rho_0+\varepsilon\rho)$ is inconsistent with (8), (24), (27) (see Remark~\ref{bpP:p11:r-cp-2}); confirmed by the review. E-P-CP09-2 (= A4): p.~21 after (47), ``$\epsilon_0\le\bar c\epsilon^{\nu_0}$'' gives only an upper bound for the admissible size; the claim needs a lower bound (Step~\ref{stepP:p11:wr-nu}, Corollary~\ref{bpP:p11a:adm-budget-beta}). E-P-CP09-3 (= A5): (27) with $\rho_0\in(\R^+)^p$, $\abs{\rho_0}=2\epsilon$ and polydisc $D^p(0,\abs{\rho_0}/2)$ may contain $\rho_j=-\rho_{0j}$ where $\sqrt{\rho_{0j}+\rho_j}$ vanishes for $p\ge2$; here all components of $\tilde\rho_0$ are equal and the polydisc has radius $\frac12\tilde\rho_{0j}$. E-P-CP09-4 (= A7): Lemma 3 for $K\cap D_{\gamma_0,\tau_0}$ (should be $K\cap\omega^{-1}(D)$), proof uses ``Theorem 2'' for Lemma 2. E-P-CP09-5 (= A6): Lemma 7 states $(\theta,r)=(\tilde x,\tilde y)$ although the slow change depends on the actions, and puts $\langle\hat\Omega,I\rangle$ twice in $Q^\star$; not used here. E-P-CP09-6 (= R-CP-2): see Remark~\ref{bpP:p11:r-cp-2}. \textbf{E-P-CP09-7 (withdrawn as a domain complaint; replaced).} The earlier claim that the domains of (16)--(17) and (28) are inconsistent is \emph{withdrawn}: restricting the pulled-back domain $D^d(0,s/(5\epsilon))\times D^p(0,1)$ after the homothety $A_\epsilon$ to the smaller sets used for the Rüssmann application is valid. What remains is a normalisation bookkeeping point: p.~10 names $\hat H_\epsilon=H_\epsilon\circ\Phi_\epsilon$ with $\Phi_\epsilon$ ``symplectic up to rescalings'' (footnote 16), whereas the final $\hat H_\epsilon$ of (28)--(30) is $\epsilon^{-1}H_\epsilon\circ(\text{chain})\circ A_\epsilon$; the conformal factor $\epsilon$ and the preserved time (a $c$-chart pulls back $H$ to $c^{-1}H\circ\Psi$, Definition~\ref{bpP:p1:chart}) are made explicit in Step~\ref{stepP:p11:wr-chain}.

\paragraph{Doubts about frozen statements.} None. The frozen P-FIXED asserts only existence; no secular non-degeneracy is a premise (Theorem~\ref{bpP:p10:interface} is proved in Part~10).

\paragraph{Forward references, open dependencies, what is received.} \emph{Proved elsewhere and cited by exact label}: (K5) is Corollary~\ref{bpP:p11a:adm-budget-beta} / Proposition~\ref{bpP:p11a:adm-budget} (File 11a-3), with $B_0=\bar\mu(\bar\mu+1)+2$ for $\rho=2$. \emph{Open dependencies on Files 11a-4/11a-5} (these labels now exist in Files 11a-4/11a-5; their proofs are cited, not re-audited here): the statement \Pref{bpP:p11a:weak-kam-quantitative} for $N=2n$, $\rho=2$, period-one angles, \emph{primary} (equivariant) $\mathcal K_b$, the injectivity and measure bound (K4) (\Pref{bpP:p11a:evaluation-bilipschitz}), the existence of the tori and Diophantine frequencies (K2) (\Pref{bpP:p11a:limit-torus}), and the displacement (K3) together with the \emph{displacement power} (K6): $\Delta\le c_D\varepsilon_f^{\kappa_D}\beta^{-a_D}$ with $c_D,\kappa_D>0$, $a_D\ge0$ depending only on the fixed entries of $\mathfrak d_*$ (\Pref{bpP:p11a:displacement}); File~11a-3 does \emph{not} supply (K6). If 11a-4/11a-5 give a weaker form (for instance $\Delta\le1/2$ for $\varepsilon_f\le\varepsilon_{\rm adm}$ only), Step~\ref{stepP:p11:wr-nu} needs just that: since $\varepsilon_f\le\varepsilon_{\rm adm}$ is already arranged, the displacement condition $\Delta\le\frac12$ is then a hypothesis on the KAM input and no power is needed; the tube inclusion uses nothing smaller. File~11a-5's \texttt{reference-constants} is not used (Lemma~\ref{bpP:p11:amount}). \emph{Rectangular averaging} (Lemma~\ref{bpP:p11:aniso-avg}) is an adaptation of Proposition~\ref{bpP:p2:nfb}; Lemma~\ref{bpP:p11:aniso-bracket} and the list (R1)--(R3) of Step~\ref{stepP:p11:aa-same} name every replaced estimate; it may be absorbed into Part~2. \emph{Part 11 supplies}: for each small $\varepsilon$ the datum $\mathfrak d_\varepsilon$ of Remark~\ref{bpP:p11:data-star}, the Hamiltonian $\hat H_\varepsilon$ with $\sup\lvert P_\varepsilon\rvert\le\varepsilon_f(\varepsilon)=O(\varepsilon^{\nu_2+3/2}\ell^{2\tau_0+1})$, and to Part~13 Theorem~\ref{bpP:p11:wrapper}. Track-N citations: \Nref{bpN:p4:norm}, \Nref{bpN:p4:sup-weighted}, \Nref{bpN:p4:coefficients}, \Nref{bpN:p4:general-cauchy}, \Nref{bpN:p4:cauchy}, \Nref{bpN:p4:algebra}; track K: \Kref{bp:p1:ext-mvt}, \Kref{bp:p1:holomorphic-analytic}.

\paragraph{Label index.} \texttt{e1}--\texttt{e5}, \texttt{kam-input}, \texttt{norm-dict}, \texttt{select}, \texttt{holo-chart}, \texttt{total-chart}, \texttt{avg-fast}, \texttt{aniso-bracket}, \texttt{aniso-avg}, \texttt{polar-avg}, \texttt{avg-slow}, \texttt{prepared}, \texttt{amount}, \texttt{r-cp-2}, \texttt{forward-tube}, \texttt{wrapper}, \texttt{data-star}, \texttt{displacement}.
\endgroup
